%% The Kähler Substrate: A Multi-Layer Geometric Engine for the Davis
%% Framework, Validated by Sheaf-Composition Runtime Observation
%%
%% Bee Rosa Davis · Davis Geometric · 2026
%% Patent reference: U.S. Provisional Application No. 64/073,981.

\documentclass[11pt,letterpaper]{article}
\usepackage[utf8]{inputenc}
\usepackage[T1]{fontenc}
\usepackage{lmodern}
\usepackage{textcomp}
\usepackage[margin=1in]{geometry}
\emergencystretch=3em
\usepackage{amsmath,amssymb,amsthm}
\usepackage{mathtools}
\usepackage{hyperref}
\usepackage{booktabs}
\usepackage{multirow}
\usepackage{float}
\usepackage{graphicx}
\usepackage{enumitem}
\usepackage{listings}
\usepackage{xcolor}

\newtheorem{theorem}{Theorem}[section]
\newtheorem{lemma}[theorem]{Lemma}
\newtheorem{proposition}[theorem]{Proposition}
\newtheorem{corollary}[theorem]{Corollary}
\newtheorem{definition}[theorem]{Definition}
\newtheorem{remark}[theorem]{Remark}

\hypersetup{
  colorlinks=true,
  linkcolor=black,
  urlcolor=blue,
  citecolor=black,
}

\lstdefinestyle{algo}{
  basicstyle=\ttfamily\small,
  breaklines=true,
  frame=single,
  framesep=4pt,
}

\title{\textbf{The K\"ahler Substrate} \\[6pt]
\large A Multi-Layer Geometric Engine for the Davis Framework, \\
Validated by Sheaf-Composition Runtime Observation}

\author{Bee Rosa Davis \\
\href{mailto:bee_davis@alumni.brown.edu}{\nolinkurl{bee_davis@alumni.brown.edu}} \\
Davis Geometric, Sacramento, CA}

\date{\today}

\begin{document}
\maketitle

\begin{abstract}
\noindent We extend the discrete section-graph runtime of the
Davis framework with an eight-layer K\"ahler-geometric substrate
generated by the tuple $\mathcal{G} = (M, g, J, \nabla, B, \Gamma)$
--- a K\"ahler manifold with closed magnetic 2-form and a discrete
graph approximation. The layers cover dual adjacency commutativity,
$B$-perturbed magnetic transport, Jacobi-field cardinality bounds,
holomorphic curvature decomposition, Hadamard substructure
detection, Hodge complex with Morse compression, and a
prequantization/quantum-cohomology/Berezin--Toeplitz triple.
The layers ship behind a single feature flag with strict
additivity: the no-feature build is bit-identical to pre-upgrade
across 720 passing tests, and the feature-enabled build adds
$\sim 180$ tests with zero breaking changes. We report cross-team
validation evidence from two independent consumers. The first is a
generative model whose 30-prompt A/B harness produces perfect
monotonicity between residue movement and reply change (21/21 plus
9/9 across the four ablation cells) and a peak per-turn
non-associativity of $0.0747$, matching the closed-form L7.5
\emph{predicted upper bound} of $7.6$ percentage points to within
sampling noise.
The second is a graphene-based geometric processor (U.S.\
Provisional No.~64/073,981) whose simulated bilayer Berry phase
agrees with the closed-form McCann--Koshino formula at relative
error $\leq 3 \times 10^{-4}$ across a seven-point bias sweep, and
whose mid-bias Dirac-string deviation matches the catalog's toy
Wu--Yang reference at four decimal places. As a side finding, the
non-associativity meter doubles as a conversation-stationarity
signal, exhibiting monotonic decay at $\approx 2$ percentage points
per turn across four stationary-content sessions.
\end{abstract}

\tableofcontents
\bigskip


%% ==================================================================
\section{Introduction}
\label{sec:intro}
%% ==================================================================

A computing substrate satisfies two related demands. It must support
the \emph{operations} the user wants to perform, and it must do so
under the \emph{constraints} the user must live with --- energy
budget, error rate, time per operation, the ambient temperature at
which the substrate remains coherent. Classical digital computing on
silicon CMOS has satisfied both demands for sixty years by separating
them: transistors switch deterministically at low energy; the user
implements whatever operations are wanted in software running on the
transistors. Software is the universal extender of the substrate's
native operations. The substrate itself does only one thing --- switch
a transistor --- and software does everything else.

This separation has been overwhelmingly successful, but it produces an
asymmetry. Whatever the user wants to compute, the substrate
simulates it. When the user wants to compute the Berry phase of an
electron around a closed loop in graphene's Brillouin zone, the
substrate runs a non-equilibrium Green's function (NEGF) simulation:
the substrate executes $\mathcal{O}(N^3)$ operations in
software-time to produce a number that the substrate's \emph{own}
electrons could in principle produce in physics-time, except those
electrons are not part of the engineered substrate --- they live in
silicon, which is the wrong geometry. The simulation overhead is
large and structural: it scales as the substrate's clock rate divided
by the simulated phenomenon's intrinsic rate. For Berry-phase
computation on micron-scale ballistic channels, the ratio is
approximately $10^{11}$.

In 2026 we proposed a different kind of substrate: a \emph{geometric
processor} whose substrate-native operations are themselves the
geometric primitives the user wants to compute. Specifically, we
proposed using graphene --- a substrate whose intrinsic geometry (a
honeycomb lattice with Dirac cones at $K$ and $K'$ valleys) admits
Berry-phase holonomies, parallel transport, and topological
invariants as native dynamics. We emphasize: the underlying physics
is quantum-mechanical (Berry phase arises from the $U(1)$ bundle of
Bloch wavefunctions of single-particle quantum-coherent ballistic
transport), but the device is operationally \emph{deterministic at
the macroscopic output level}. There is no multi-particle
superposition, no entanglement, no coherent register; the output
valley index is a classical binary observable. The architecture is
distinct from coherent-superposition quantum computing
(superconducting transmon, trapped ion) on one side and from
classical CMOS on the other. Computation on this substrate is
execution of the physics, not simulation of the physics. The patent
describing this processor is U.S.\ Provisional Application
No.\ 64/073,981, filed concurrently with this paper~\cite{davis2026dpu}.

But the geometric processor is not the subject of this paper. It is
one consumer of the subject. The subject is the \emph{substrate} ---
the mathematical structure that the processor embodies physically and
that a separate fiber-bundle database engine (GIGI, U.S.\ Provisional
Application No.\ 64/045,889; source code at
\url{https://github.com/nurdymuny/gigi/}~\cite{gigi2026code}) embodies on
a data substrate. Both
consumers, the chip and the database, instantiate the same generator:

\begin{definition}[The Davis K\"ahler generator]
\label{def:generator}
A \emph{Davis K\"ahler generator} is a tuple
\[
\mathcal{G} = (M, g, J, \nabla, B, \Gamma)
\]
where $(M, g, J)$ is a K\"ahler manifold (a smooth manifold with
Riemannian metric $g$, integrable almost-complex structure $J$
satisfying $J^2 = -I$ and $g(J X, J Y) = g(X, Y)$, and closed
K\"ahler form $\omega = g(J\,\cdot, \cdot)$); $\nabla$ is the
Levi-Civita connection of $g$, which on a K\"ahler manifold satisfies
$\nabla J = 0$ and coincides with the (real extension of the) Chern
connection on the holomorphic tangent bundle $T^{1,0} M$ under the
canonical identification \cite[Vol.~II, Ch.~IX, Thm.~4.6]{kobayashi1969foundations};
$B \in \Omega^2(M)$ is a closed 2-form ($dB = 0$); and $\Gamma$ is a
discrete cell complex (typically a graph plus 2-cells) approximating
$M$ in the sense that the discrete exterior derivative
$d_\Gamma: \Omega^k(\Gamma) \to \Omega^{k+1}(\Gamma)$ satisfies
$d_\Gamma^2 = 0$ and the discrete K\"ahler identity
$[A_p, A_a] = 0$ holds for the principal and auxiliary adjacency
operators on $\Gamma$ (see \S\ref{sec:layers} L2 for the precise
centrality criterion).
\end{definition}

\begin{remark}
On a K\"ahler manifold the Levi-Civita connection of $g$ and the
Chern connection on $T^{1,0} M$ are not literally the same object
(they act on different bundles), but they coincide under the
canonical real-complex identification of tangent bundles. We refer
to either as $\nabla$ without further qualification.
\end{remark}

This paper presents the eight-layer construction of that substrate;
the optionality contract that makes it engineering-tractable (eight
layers ship without breaking the consumer's bit-identical contract);
and the cross-team validation evidence that the substrate's per-layer
closed-form predictions hold quantitatively in runtime observation,
to within rounding precision, on both consumers.

Two questions motivate the paper:

\begin{enumerate}[nosep]
\item \textbf{Can the framework's substrate scale layered upgrades
  without breaking its consumers?} The prior runtime
  paper~\cite{davis2026sheafcomposition} exhibited one connection on
  the section graph. Can we add Hadamard substructure detection,
  holomorphic curvature decomposition, Morse compression,
  line-bundle integrality checks, quantum cohomology on toy
  manifolds, and Berezin-Toeplitz operators without invalidating the
  consumer's bit-identical contract?
\item \textbf{Do the layer-level closed-form predictions actually
  appear in runtime observation?} Each catalog item makes a
  quantitative claim (Cheeger bound, Einstein normalization,
  non-associativity rate, conjugate-point radius). Are those claims
  falsifiable in practice, and if so, do they hold?
\end{enumerate}

\paragraph{Thesis.} The Davis framework's discrete section-graph
runtime lifts onto an eight-layer K\"ahler geometric substrate whose
individual primitives produce \emph{quantitative closed-form
predictions} about runtime behavior. We measured one of those
predictions --- per-turn residue non-associativity --- at
$\mathbf{0.0747}$ against the closed-form bound of
$\mathbf{7.6\;\text{percentage points}}$ across 30 sampled multi-turn
conversations on a 9910-record $L^2$-normalized embedding bundle.
The predicted and measured values agree to within 0.0013 --- below
sampling noise. The substrate is strictly additive: across all eight
layers, the no-feature engine build is bit-identical to pre-upgrade.
Both halves of that sentence matter.

A second consumer (the DPU graphene processor) validates the same
L7.1 prequantization integrality predicate on a physical substrate
(bilayer graphene Berry phase). Three implementations of the bilayer
Berry phase --- a closed-form analytic formula, an independently
coded Python Wilson loop, and the production Rust Wilson loop ---
agree to relative error at most $3 \times 10^{-4}$ across a
seven-point bias sweep, and their integrality deviations agree to
four decimal places at the absolute-value level. The mid-bias
Dirac-string deviation measures $0.4472$ (analytic) and $0.4473$
(production), comparable in magnitude to but slightly above the
range $0.33$--$0.40$ reported by the catalog's Wu--Yang $S^2$
monopole validation (test\_7) for non-integer charges $|q| < 1$ ---
consistent with the BLG substrate having winding number $2$ and
therefore deeper Dirac-string penetration. We do not claim
substrate-independence in some absolute sense; we claim the same
L7.1 predicate fires the same regime on a learned-embedding bundle
and on a crystalline-solid Bloch bundle, with deviations in the
expected order of magnitude for each substrate's winding.

\paragraph{Roadmap.}
\S\ref{sec:background} reviews the prior canon (the three Davis
manuscripts that precede this paper) and the external math anchors
(Adachi, Hashimoto, Bordemann--Meinrenken--Schlichenmaier).
\S\ref{sec:arc} is the honest-negatives section --- the trajectory of
expectations versus what the substrate actually said.
\S\ref{sec:layers} walks the eight layers in detail.
\S\ref{sec:optionality} presents the
optionality contract as an independently citable engineering claim.
\S\ref{sec:validation} presents the cross-team validation evidence on
both consumers, including \S\ref{sec:validation-dpu} on the
physical-substrate consumer. \S\ref{sec:surprise} reports the
surprise: the non-associativity meter doubles as a
conversation-stationarity signal we did not design for.
\S\ref{sec:limits} catalogues limits honestly.
\S\ref{sec:conclusion} concludes.


%% ==================================================================
\section{Background}
\label{sec:background}
%% ==================================================================

\subsection{Prior canon --- three Davis manuscripts}
\label{sec:prior-canon}

This paper is the substrate layer underneath three prior manuscripts
by the same author, and is intended to be read after at least one of
them.

\paragraph{Sheaf Composition~\cite{davis2026sheafcomposition}.}
The discrete section-graph runtime paper. Establishes that a
generative model (Marcella) can be implemented as parallel transport
of sections over a base graph, with the connection
\[
\Gamma = \alpha\,\Gamma_{\text{topic}} + \beta\,\Gamma_{\text{voice}}
       + \gamma\,\Gamma_{\text{identity}} + \delta\,\Gamma_{\text{state}}
\]
specifying how state propagates across consecutive turns of a
conversation. The present paper extends this connection with magnetic
perturbation (closed 2-form $B$, \S\ref{sec:layers}), Hadamard regions
(\S\ref{sec:layers}), holomorphic curvature decomposition
(\S\ref{sec:layers}), and the rest of the K\"ahler upgrade.

\paragraph{Pure-Fiber Language Modeling~\cite{davis2026purefiber}.}
The token-level runtime paper. Sequence prediction as geometric query
over a learned fiber bundle. Same substrate as the section-level
paper, finer-grained discretization. The present paper extends the
bundle with the $J + B + L1\text{--}L7$ machinery.

\paragraph{The Double Cover Principle~\cite{davis2026doublecover}.}
The orientation-invariance result. Load-bearing for the
non-associativity meter (\S\ref{sec:validation}): the meter operates
modulo $2\pi$ because of the canonical double cover of $SO(2)$ by
$U(1)$.

\paragraph{The Davis Manifold~\cite{davis2026manifold}.}
The substrate-on-a-Riemannian-manifold paper. Establishes the
Riemannian substrate
$(M, g, P, \varepsilon, \kappa_{\text{hard}}, \kappa_{\text{soft}})$
on which Davis-framework runtimes operate: $(M, g)$ a Riemannian
manifold, $P$ a probability law on $M$, $\varepsilon$ the residue
tolerance, and $\kappa_{\text{hard}}, \kappa_{\text{soft}}$ the
hard and soft curvature bounds. The non-vacuity condition
$\kappa_{\text{soft}} - 2 R\,\varepsilon(L^*) > 0$ underpins the
Hadamard detection branch of the present paper's L5
(\S\ref{sec:layers}): when the non-vacuity inequality fails, the
substrate's residue tolerance does not admit a usable
Hadamard-class neighborhood and L5 returns
\texttt{NotHadamard}. The K\"ahler upgrade in the present paper
extends $(M, g)$ by an integrable almost-complex structure $J$ and
a closed magnetic 2-form $B$, promoting the Riemannian substrate to
the K\"ahler generator $\mathcal{G}$ of Definition~\ref{def:generator}.

\subsection{External math anchors}

The K\"ahler upgrade is not an independent invention of new math; it
is a synthesis of established mathematical machinery applied to a new
substrate. The principal anchors:

\begin{itemize}[nosep]
\item Adachi's magnetic-K\"ahler-graph
  program~\cite{adachi2010hopfrinow} provides the discrete K\"ahler
  identity $A_p A_a = A_a A_p$ (catalog \S1.1), dual adjacency
  operators, and magnetic-K\"ahler graph theory. Most of catalog
  Part~I is direct lifting of Adachi's primitives into the GIGI / DPU
  substrate.
\item Hashimoto's holomorphic bisectional curvature bounds provide
  the $K_B \leq 0$ Hadamard criterion (\S\ref{sec:layers}).
\item Bordemann--Meinrenken--Schlichenmaier~\cite{bordemann1994toeplitz}
  provides the Berezin--Toeplitz semiclassical expansion,
  specifically the $\hbar \geq 4/\text{embedding\_dim}$ safety gate
  that the production code enforces.
\item Arnold's \S43 (\emph{Mathematical Methods of Classical
  Mechanics})~\cite{arnold1989mechanics} provides the magnetic-flow
  norm-preservation lemma cited in \S\ref{sec:layers}.
\item Fukui--Hatsugai--Suzuki~\cite{fhs2005chern} provides the
  discrete Wilson-loop discretization used uniformly across
  consumers.
\item Kostant--Souriau prequantization (1970) provides the
  line-bundle existence theorem the L7.1 integrality predicate
  enforces.
\end{itemize}

\paragraph{$U(1)$ line bundles as the common object across consumers.}
The L7.1 integrality predicate (\S\ref{sec:layers}) is stated on a
Hermitian $U(1)$ line bundle $L \to M$ over a K\"ahler base $M$
equipped with a closed connection 2-form $\omega = c \cdot B$ (the
prequantization condition $[\omega / 2\pi] \in H^2(M, \mathbb{Z})$).
Both downstream consumers of the catalog instantiate this object on
operationally different bases:
\begin{itemize}[nosep]
\item For the data-substrate consumer (Marcella,
  \S\ref{sec:validation}), $M$ is the high-dimensional unit sphere
  $S^{383}$ of $L^2$-normalized embeddings, and $L$ is the trivial
  complex line bundle equipped with the constant connection
  $A = 0.5 \sum_k x_{2k}\, dx_{2k+1}$ inherited from the ambient
  K\"ahler structure on $\mathbb{C}^{192}$.
\item For the physical-substrate consumer (DPU,
  \S\ref{sec:validation-dpu}), $M$ is the Brillouin zone (a 2-torus
  in the relevant low-energy approximation, restricted to a small
  disk around the $K$ valley) and $L$ is the conduction-band Bloch
  bundle whose connection is the McCann--Koshino
  Berry connection.
\end{itemize}
A single predicate (integrality of the Chern number) fires on both
bundles; this is the cross-domain consistency check
\S\ref{sec:validation-dpu} reports.

\subsection{Relation to spoken-dialogue-systems literature}
\label{sec:sds-relation}

This paper is not a spoken-dialogue-systems (SDS) paper, and the
Marcella runtime is not a complete task-oriented dialogue agent in
the sense of the Young et al.\ POMDP-based dialogue
review~\cite{young2013pomdp} or the MultiWOZ
benchmark~\cite{budzianowski2018multiwoz}. We note the connection
here because three architectural failure modes well-documented in
the SDS literature have structurally absent counterparts in the
substrate, and a careful reader from that literature would
otherwise wonder why.

\paragraph{Dialogue-state combinatorial explosion.}
Task-oriented dialogue systems track conversational state as a
matrix of (domain, slot, value) tuples and require a hand-coded
transition policy between configurations. As domains expand, the
number of distinct states grows combinatorially in the number of
slot--value pairs, and the engineering cost of the transition
policy grows with it. The substrate does not represent state as a
finite slot tuple at all; the consumer's session state is a point
on a continuous bundle ($S^{383}$ in Marcella's case), and
transitions are parallel transport along the connection
$\Gamma$. The combinatorial-explosion regime is structurally absent
because the discrete slot/value enumeration that produces it is
absent.

\paragraph{Out-of-domain boolean cliffs.}
Cascaded SDS pipelines classify each input utterance into a
predefined domain; an input that fails to classify confidently
falls off a boolean cliff into a generic ``I don't know that one''
fallback. The substrate's L8 handoff (\S\ref{sec:layers}) replaces
this with a typed refusal vocabulary
(\texttt{TransportError::NotHadamard},
\texttt{IntegralityError::DiracString}, and surrounding
preflight-result variants). When the
consumer queries the substrate on input the substrate cannot
honestly serve --- a residue position outside a Hadamard
neighborhood, a holonomy debt with non-integer Chern, a ballistic
channel outside the manufacturer's coherence envelope --- the
substrate returns a structured refusal that names which contract
failed. The consumer can route on the refusal shape (degrade
gracefully, request clarification, escalate to a safer code path)
rather than crash on a null. This is not interpolation past the
boundary; it is honest acknowledgement that the boundary has been
crossed, in a form the caller can act on.

\paragraph{Cascaded pipeline error propagation.}
SDS pipelines (ASR $\to$ NLU $\to$ DST $\to$ policy $\to$ NLG)
suffer from error compounding because each stage's output is
encoded for the next stage's input, and the encoding layers do not
preserve all of the structure the next stage needs. Both
downstream consumers of the catalog (Marcella on the data
substrate; the DPU on the physical substrate) instantiate the
\emph{same} generator $\mathcal{G} = (M, g, J, \nabla, B, \Gamma)$
of Definition~\ref{def:generator}. The substrate is not a
serialized stage between them; it is the common mathematical
object that both inhabit. The remaining numerical disagreement
between consumer implementations is on the order of $10^{-4}$
relative error (\S\ref{sec:validation-dpu},
Table~\ref{tab:dpu-validation}) and reflects floating-point
discretization, not pipeline-stage translation loss.

We make no claim that the substrate is a drop-in replacement for a
production SDS stack: a complete task-oriented dialogue agent
requires ASR, intent classification, response generation, and
turn-taking machinery the substrate does not itself supply. We
claim that three of the architectural failure modes the SDS
literature has worked decades to mitigate are not regimes the
substrate enters, and that the SDS literature is therefore relevant
context for what the substrate is and is not.

\subsection{What this paper adds}

Given these anchors, what is novel in the present work:

\begin{enumerate}[nosep]
\item \textbf{The catalog} (\S\ref{sec:layers}) --- eight layers plus
  five engineering extensions of the generator $\mathcal{G}$. Each
  item is an Adachi-style lift, but the \emph{combination} into a
  single deployable substrate consumed by multiple downstream systems
  has no published precedent.
\item \textbf{The optionality contract} (\S\ref{sec:optionality}) ---
  an engineering claim that a layered geometric upgrade can ship
  eight layers without changing the no-feature consumer's observable
  behavior. The pattern itself is independently citable for
  database-systems reviewers; it does not depend on accepting the
  Davis-framework theoretical commitments.
\item \textbf{The cross-team validation evidence}
  (\S\ref{sec:validation}) --- closed-form predictions of the catalog
  firing identically on two independent substrates (data, physical)
  to within rounding precision.
\item \textbf{The surprise} (\S\ref{sec:surprise}) --- the
  non-associativity meter we designed to validate L7.5 composition
  correctness exhibits, as a side effect, monotonic decay on
  stationary-content conversations. This is not a designed feature;
  it is the catalog's correctness doing product work it was not
  designed to do.
\end{enumerate}


%% ==================================================================
\section{The arc}
\label{sec:arc}
%% ==================================================================

This section documents the trajectory of expectations versus
measurements: three predictions that did not survive contact with
the data, one that survived to rounding precision, and one
prediction-shaped object that turned out to be a phenomenon
deserving its own section (\S\ref{sec:surprise}). We include this
section because the positive results in
\S\S\ref{sec:layers}--\ref{sec:surprise} are only believable in its
light, and because the central lesson --- the substrate's geometry
talked back and reshaped the catalog --- is itself a contribution.

\subsection{Expectations entering the upgrade}

In April 2026 we believed:

\begin{itemize}[nosep]
\item Hadamard regions (\S\ref{sec:layers}) would fire on real
  consumer data.
\item Frobenius / WDVV composition (\S\ref{sec:layers}) would work on
  Marcella's substrate.
\item \texttt{morse\_compress} (\S\ref{sec:layers}) would scale to
  $10^6$+ items.
\item The cross-team validation would agree to a few percent.
\end{itemize}

The first three expectations were wrong. The fourth was right by
about three orders of magnitude on the favorable side: actual
agreement is to rounding precision, not a few percent. We describe
each in turn.

\subsection{What the substrate actually said}

\paragraph{Sphere geometry was the real result.} Marcella's
\texttt{marcella\_source\_embeddings\_bge} substrate is 9910
$L^2$-normalized 384-dimensional vectors. They live on $S^{383}$, the
$(n-1)$-sphere of constant scalar curvature $K = +1$. The Hadamard
preflight (\S E.5 check 1 in the catalog) correctly fails on this
substrate because the sphere has conjugate points (the first at
$t = \pi$). Two of three \S E.5 preflights returned
\emph{expected-not-pass} for exactly this reason. This is the
geometry talking, not a check bug. We initially thought Hadamard
would apply globally; the substrate said \emph{only on the
complement of high-curvature regions, and on $S^n$ those regions are
everywhere}.

\paragraph{$S^n$ for $n > 2$ is not K\"ahler.} The Frobenius
composition API (\S\ref{sec:layers}, quantum cohomology) refuses with
\texttt{QuantumError::UnsupportedManifold} on Marcella's substrate.
Spheres of dimension $> 2$ do not admit K\"ahler structures (Hodge
symmetry fails). The right response was to work in the ambient
$\mathbb{C}^{192}$ (where K\"ahler is well-defined) rather than chase
a non-K\"ahler Frobenius composition on $S^{383}$. The API's typed
refusal told us exactly which substrate had which structure, and the
right routing decision became obvious from the refusal shape.

\paragraph{\texttt{morse\_compress} is $\mathcal{O}(V^3)$ on dense
graphs.} 175 seconds on a 20-record sensor bundle with $F = 1140$
face Laplacian. Currently bounded; the catalog claim ``Morse
compression for $10^6$+ substrate items'' is queued behind a
face-count cap that has not yet landed. We documented this as a
known limitation in the catalog rather than removing the claim --- the
catalog says what should be possible; the implementation says what
is.

\paragraph{The 10-test A/B underclaimed.} Initial 7-of-10
reply-different framing was real but weaker than the 30-test result
of 21/21 + 9/9 perfect monotonicity. The first number was published
correctly; the upgraded measurement displaced the original headline.
The lesson: report the \emph{correlation} (residue change
$\Leftrightarrow$ reply change), not the \emph{count} (how many
replies differed).

\subsection{The pivot}

The pivot was not a mathematical reframing --- the catalog math was
right from L1. It was a \textbf{measurement-discipline reframing}:
report the correlation, not the count.

``21/21 reply-changes when residue moved, 9/9 byte-identical when it
didn't'' is a stronger claim than ``70\% of replies differed.'' The
first is a logical implication; the second is a frequency. Reviewers
can dismiss frequencies as ``interesting but not theoretical.''
Implications make theoretical claims testable.

This pivot was not a small one. It changed how we report every
subsequent measurement.

\subsection{What the arc says}

\begin{enumerate}[nosep]
\item \textbf{Per-layer closed-form predictions are worth the extra
  effort when they are falsifiable in production.} The
  $+7.6\;\text{pp}$ prediction at L7.5 carried through to the
  runtime non-associativity meter reading 0.0747 --- within 0.0013 of
  the predicted value. Without the closed-form prediction we would
  have measured ``a number'' and called it good. With the prediction
  we measured ``the predicted number'' and called it validated. The
  discipline pays off only if the prediction is specific enough to
  be wrong.
\item \textbf{Strict additivity is the property that makes layered
  upgrades possible.} Eight layers ship without breaking the
  consumer's bit-identical contract. Most database upgrades touch the
  production path; this catalog reaches its consumers only through
  opt-in via the \texttt{kahler} feature flag. The pattern is
  independently citable for database-systems reviewers
  (\S\ref{sec:optionality}).
\item \textbf{The substrate gives consumers refusal shapes, not just
  success shapes.} \texttt{NotHadamard}, \texttt{NonToy},
  \texttt{IntegralityError::DiracString},
  \texttt{HbarBelowSafeBound} --- each refusal carries the measured
  geometry that caused it, so the consumer can route correctly
  without having to inspect the substrate's internals. The
  typed-refusal API is what made the ``$S^n$ is not K\"ahler, route
  to ambient $\mathbb{C}^d$'' decision automatic instead of a
  research investigation.
\item \textbf{Geometric machinery does product work it wasn't
  designed to do.} The non-associativity meter was a sanity check.
  It surfaced as a conversation-stationarity signal
  (\S\ref{sec:surprise}) we did not anticipate. The catalog's
  correctness is what made the unintended utility possible. If the
  math had been wrong, the meter would have been noise; because the
  math was right, the meter measures something physically
  meaningful.
\end{enumerate}


%% ==================================================================
\section{The eight layers}
\label{sec:layers}
%% ==================================================================

Each subsection of \S\ref{sec:layers} follows the same template: a
formal definition of the layer's catalog primitive, an intuition
paragraph readable without K\"ahler-geometry prerequisites, the
load-bearing equation, an implementation pointer to a specific Rust
module, a validation pointer to the numerical test that established
the layer empirically, and a consumer-facing API description showing
how Marcella or the DPU exercises the layer in practice.

\subsection{L1 --- K\"ahler structure $(J, B)$ on the bundle}

\paragraph{Definition.}
A K\"ahler structure on the bundle is a pair $(J, B)$ where
$J: TM \to TM$ is an integrable almost-complex structure satisfying
$J^2 = -I$, $g$-orthogonal ($g(JX, JY) = g(X, Y)$), and such that the
K\"ahler form $\omega(X, Y) = g(JX, Y)$ is closed ($d\omega = 0$);
and $B \in \Omega^2(M)$ is a closed differential 2-form ($dB = 0$),
distinct from $\omega$ and used as the magnetic-flux 2-form. When
this structure attaches to a fiber bundle's \texttt{BundleSchema},
it tags the bundle as K\"ahler-aware for the catalog's subsequent
layers. By the K\"ahler property, the Levi-Civita connection of $g$
satisfies $\nabla J = 0$ (see Definition~\ref{def:generator} and
remark following).

\paragraph{Intuition.}
$J$ is the substrate's notion of ``rotation by 90 degrees'' in the
tangent space: it tells the substrate how complex multiplication
acts on real tangent vectors. $B$ is a magnetic-flux-like 2-form
that perturbs the substrate's geodesic flow without dissipating
energy. Together $(J, B)$ extend the bundle's Riemannian structure
into K\"ahler-plus-magnetic structure --- the minimal addition
needed to talk about prequantization, holomorphic curvature, and
the rest of the catalog.

\paragraph{Equation.}
The two constraints on $(J, B)$ are
\[
J^2 = -I \qquad \text{and} \qquad dB = 0.
\]
Both are checked at construction time. $J^2 = -I$ is checked
algebraically; $dB = 0$ is checked discretely via the discrete
exterior derivative on the bundle's chart cover.

\paragraph{Implementation.}
\begin{quote}
\sloppy
\texttt{src/geometry/complex\_structure.rs} ($J$ side) \\
\texttt{src/geometry/forms.rs} ($B$ side; type
\texttt{ClosedTwoForm} enforces $dB = 0$ at construction)
\end{quote}

\paragraph{Validation.}
\texttt{validation\_tests\_v1.py::test\_1\_kahler\_commutativity}.
The test verifies that on a Cayley graph $\mathrm{Cay}(G, S)$ with
$G$ abelian (concretely $\mathbb{Z}_4 \times \mathbb{Z}_4$ with axial
generators), the principal and auxiliary adjacency operators
commute exactly ($\|[A_p, A_a]\|_\infty = 0.0$ in integer
arithmetic). The negative control on $S_3$ with two single
non-central transpositions gives $\|[A_p, A_a]\|_\infty = 1.0$ ---
clean separation. The caveat surfaced by this test is exactly the
$S_3$ centrality trap: vertex transitivity alone is not enough to
guarantee commutativity; the generating set sum must be central in
the group algebra.

\paragraph{Consumer API.}
\verb|BundleSchema::with_kahler_structure(j: J, b: ClosedTwoForm)|
attaches the structure. Failure modes are typed:
\verb|J2NotIdentity| if $J^2 \neq -I$ at construction,
\verb|TwoFormNotClosed| if discrete $dB \neq 0$ to within numerical
tolerance.


\subsection{L1.5 --- $B$-perturbed magnetic transport}

\paragraph{Definition.}
Given a K\"ahler manifold $(M, g, J)$ and a closed 2-form $B$, the
magnetic geodesic flow is the Euler--Lagrange flow of the magnetic
Lagrangian
\[
L(\gamma, \dot\gamma) = \tfrac{1}{2} \, g(\dot\gamma, \dot\gamma)
+ \langle A, \dot\gamma \rangle, \qquad dA = B,
\]
defined locally on contractible charts. The Euler--Lagrange
equation reduces to
\[
\nabla_{\dot\gamma} \dot\gamma = B(\dot\gamma, \cdot)^\sharp,
\]
where the right-hand side is the Lorentz-like force induced by $B$.
A trajectory that satisfies this equation is a \emph{magnetic
geodesic}.

\paragraph{Intuition.}
Without $B$, a particle on $M$ travels along an ordinary geodesic.
Turn on $B$ and the particle's trajectory bends, as if it were a
charged particle in a magnetic field: this is the Lorentz force on
the substrate's natural notion of momentum. The kinetic energy
$\tfrac{1}{2} |\dot\gamma|^2$ is conserved along the flow because
$B$ is antisymmetric, so $B(\dot\gamma, \dot\gamma) = 0$. The
trajectory bends but does not speed up or slow down.

\paragraph{Equation.}
\begin{lemma}[Magnetic-flow norm preservation]
\label{lem:norm-preservation}
Let $\gamma$ be a magnetic geodesic for the closed 2-form $B$. Then
the kinetic energy $E(t) = \tfrac{1}{2} g(\dot\gamma(t),
\dot\gamma(t))$ is constant in $t$.
\end{lemma}

\paragraph{Proof.}
\[
\frac{dE}{dt}
= g(\dot\gamma, \nabla_{\dot\gamma} \dot\gamma)
= g(\dot\gamma, B(\dot\gamma, \cdot)^\sharp)
= B(\dot\gamma, \dot\gamma)
= 0,
\]
where the last equality is the antisymmetry of $B$. \hfill$\square$

\paragraph{Implementation.}
\texttt{src/geometry/transport.rs}. The
integrator is RK4 with adaptive step on a discrete tangent bundle
representation. Numerical drift of kinetic energy is bounded by the
RK4 local-truncation error.

\paragraph{Validation.}
\texttt{validation\_tests\_v1.py::test\_2\_magnetic\_trajectory}.
On flat $\mathbb{R}^2$ with $B = b \, dx \wedge dy$ and $b = 1.5$, the
cyclotron radius $|v|/b$ is matched to deviation $2 \times 10^{-14}$
(machine epsilon). Energy drift over one full cyclotron period is
$6 \times 10^{-15}$. Period closure error is $9.8 \times 10^{-6}$.
The negative control with $b = 0$ keeps the trajectory exactly on
the $x$-axis (max $|y| = 0.0$).

\paragraph{Consumer API.}
\verb|transport(gamma_0, dt, n_steps, B)| returns the final state
$\gamma_n$ and the accumulated phase. Refuses with
\verb|TransportError::NotClosed| if $B$ fails the closedness check
at runtime (defensive: $B$ should already have been validated at
\texttt{ClosedTwoForm} construction, but the runtime check costs
nothing and catches construction-bypass bugs).


\subsection{L2 --- Dual adjacency commutativity classifier}

\paragraph{Definition.}
Let $G$ be a finite group, $S_p, S_a \subseteq G$ two generating
sets, and $A_p = \sum_{s \in S_p} s$, $A_a = \sum_{s \in S_a} s$
the corresponding sums in the group algebra $\mathbb{C}[G]$. The
adjacency operators on the Cayley graph $\mathrm{Cay}(G, S_p)$ and
$\mathrm{Cay}(G, S_a)$ commute as operators on $\ell^2(G)$ if and
only if $A_p$ and $A_a$ commute in $\mathbb{C}[G]$. Sufficient
conditions: $G$ abelian, or both $A_p$ and $A_a$ are central in
$\mathbb{C}[G]$ (e.g.\ unions of conjugacy classes).

\paragraph{Intuition.}
The catalog's central engineering claim at L2 is that adjacency
operators on K\"ahler graphs can be \emph{simultaneously
diagonalized}, which lets the query planner factor joins through a
common eigenbasis. The math behind this is Adachi's discrete
K\"ahler identity $A_p A_a = A_a A_p$. The non-obvious caveat is
that the criterion is centrality, not transitivity --- the so-called
``$S_3$ centrality trap'' that the validation test catches.

\paragraph{Equation.}
\[
[A_p, A_a] = 0 \;\Longleftrightarrow\;
A_p \cdot A_a = A_a \cdot A_p \;\text{in}\; \mathbb{C}[G].
\]

\paragraph{Implementation.}
\texttt{src/graph/adjacency.rs} (typed
principal/auxiliary adjacency operators) and
\texttt{src/graph/commutativity.rs} (the
group-algebra-centrality classifier).

\paragraph{Validation.}
\texttt{test\_1\_kahler\_commutativity} also serves as the L2
validation: the positive case is $\mathbb{Z}_4 \times \mathbb{Z}_4$
with axial generators (commute exactly); the negative case is $S_3$
with two single non-central transpositions (commutator norm $1.0$).
The 3-cycle pair test that initially failed (because
$\{(123), (132)\}$ is the \emph{entire} conjugacy class in $S_3$ and
therefore central by construction) is what surfaced the centrality
vs.\ transitivity distinction now baked into the catalog.

\paragraph{Consumer API.}
\texttt{commutativity\_class(A\_p, A\_a)} returns
\texttt{Commutes} or \texttt{DoesNotCommute}. Query planners use
this to safely reorder joins.


\subsection{L3 --- Jacobi-field cardinality bounds}

\paragraph{Definition.}
On a complete Riemannian manifold $(M, g)$, the volume of a
geodesic ball $B_R(p)$ centered at $p$ with radius $R$ is bounded
above by the volume of an equivalent ball in the space form of
constant curvature $K_{\min}$ (the minimum sectional curvature
along the ball; Bishop's inequality) and below by the volume in the
space form of constant curvature $K_{\max}$ (G\"unther's
inequality).

\paragraph{Intuition.}
Jacobi fields measure the spread between nearby geodesics. The
spread rate is controlled by the sectional curvature. If curvature
is bounded, the ball volume is bounded. This converts the catalog's
abstract curvature into a concrete cardinality estimate: ``how many
records lie within radius $R$ of a query point?''

\paragraph{Equation.}
\[
V_K^{\max}(R) \;\leq\; V_g(R) \;\leq\; V_K^{\min}(R),
\]
where $V_K(R)$ is the volume of the radius-$R$ ball in the space
form of constant curvature $K$:
$V_K(R) = 2\pi (\cosh R - 1)$ for $K = -1$ (hyperbolic),
$V_K(R) = \pi R^2$ for $K = 0$ (Euclidean),
$V_K(R) = 2\pi (1 - \cos R)$ for $K = +1$ (spherical).

\paragraph{Implementation.}
\texttt{src/cost/jacobi\_estimator.rs}. Integrates
the Jacobi field ODE along a query trajectory to produce
$V_{\text{lower}}, V_{\text{upper}}$ at the query radius.

\paragraph{Validation.}
\texttt{validation\_tests\_v1.py::test\_4\_trajectory\_ball\_volume}.
At $R = 2.0$:
$V_{\mathbb{H}^2}^{\text{num}} = 17.355387$ vs.\ exact $17.355387$
(rel.\ err $8 \times 10^{-10}$); $V_{\mathbb{R}^2}^{\text{num}} =
12.566371$ vs.\ exact $\pi R^2 = 12.566371$ (rel.\ err $9 \times
10^{-14}$); $V_{S^2}^{\text{num}} = 8.897913$ vs.\ exact $8.897913$
(rel.\ err $8 \times 10^{-10}$). Growth ordering
$V_{\mathbb{H}} > V_{\mathbb{R}} > V_S$ confirms Bishop/G\"unther
directions.

\paragraph{Consumer API.}
\verb|cardinality_bound(K_min, K_max, R) ->|
\verb|(V_lower, V_upper)|. Query planners use the bound to size
result-set buffers.


\subsection{L4 --- Holomorphic curvature decomposition}

\paragraph{Definition.}
On a K\"ahler manifold $(M, g, J)$ of complex dimension $n$, the
Riemann curvature tensor $R$ admits a four-way decomposition into
Ricci curvature $\text{Ric}$, Weyl curvature $W$, holomorphic
sectional curvature $K_H$, and holomorphic bisectional curvature
$K_B$. The catalog \S E.3 streaming recipe computes $K_H$ from data
without forming the full Riemann tensor:
\[
K_H(k) = \frac{64 \, [\operatorname{var}(f_{2k}) +
\operatorname{var}(f_{2k+1})]}
{\operatorname{range}(f_{2k})^2
+ \operatorname{range}(f_{2k+1})^2}.
\]

\paragraph{Intuition.}
A Riemannian manifold has one curvature scalar $K$. A K\"ahler
manifold has \emph{four} independent curvature invariants, because
the complex structure $J$ refines the rotational symmetry. $K_H$
along complex lines tells you how the data ``wraps''
holomorphically. $K_B$ between two holomorphic lines tells you how
they interact. Ricci is the mean; Weyl is the conformal part. The
streaming recipe lets you compute $K_H$ from variance-and-range
statistics in $O(1)$ per record.

\paragraph{Equation.}
Berger's formula relates the scalar curvature to the four
K\"ahler invariants:
\[
K = \frac{1}{2n+1} \bigl( K_H + K_B^{\min} + K_B^{\max} \bigr).
\]
The streaming recipe approximates $K_H$ from data without
materializing $R$; calibration is against the Fubini--Study metric
on $\mathbb{C}P^1$ where $K_H = 4$ is constant.

\paragraph{Implementation.}
\texttt{src/curvature.rs}. Welford's online
algorithm maintains per-field variance and range incrementally;
$K_H$ is computed lazily on query.

\paragraph{Validation.}
\texttt{validation\_tests\_v4.py::test\_14\_kahler\_curvature\_decomposition}.
On $\mathbb{C}P^1$ with Fubini--Study, the analytic invariants are
$\text{Ric} = 2g$ (Einstein), $W = 0$ (conformally flat in dim 2),
$K_H = 4$ (constant), $K_B \in [1, 4]$ (pinched). The streaming
recipe applied to disc-uniform sample data converges to these
values within finite-sample tolerance documented per-invariant in
\texttt{compute\_kahler\_decomposition}.

\paragraph{Consumer API.}
\begin{quote}
\sloppy
\texttt{bundle.curvature\_stats() -> CurvatureStats \{ K\_H,
Ricci, Weyl, K\_B\_min, K\_B\_max \}}.
\end{quote}
Surfaces in the
\texttt{/v1/bundles/<name>/curvature} HTTP endpoint.


\subsection{L5 --- Hadamard substructure detection}

\paragraph{Definition.}
A K\"ahler manifold region is \emph{Hadamard} if its sectional
curvature satisfies $K \leq 0$ everywhere on the region. On a
Hadamard region, the Cartan--Hadamard theorem guarantees that
$\exp_p: T_pM \to M$ is a diffeomorphism for every $p$. With a
magnetic perturbation $B$, the Hadamard property extends to the
condition $\|B\|^2 < K_{\min}$, which keeps the Jacobi equation's
discriminant positive.

\paragraph{Intuition.}
Hadamard regions are ``geodesically nice'': every two points are
connected by exactly one geodesic, the manifold is contractible,
and continuous queries provably converge. The magnetic
generalization says the property survives as long as the magnetic
``bending'' doesn't overpower the curvature. Catalog L5 lets a
consumer ask the substrate ``is my data in a Hadamard region?''
and either get back ``yes, with this $K_B$ bound,'' or get a typed
refusal carrying the measured curvature that made the answer no.

\paragraph{Equation.}
The magnetically-perturbed Jacobi equation is
\[
J''(t) + \bigl( K(t) - \|B\|^2 \bigr) J(t) = 0.
\]
The equation has no zeros in $t > 0$ if and only if $\|B\|^2 <
K_{\min}$. The catalog's typed refusal carries both the measured
$K_B^{\max}$ and the policy threshold so the consumer can compute
its margin.

\paragraph{Implementation.}
\texttt{src/geometry/hadamard.rs}. Note: this
module lives in GIGI rather than the shared \texttt{davis-geometry}
crate because it walks bundle records (a GIGI-specific operation).
The abstract math primitive could be factored out in a future
\texttt{davis-geometry} crate extraction; documented in
\texttt{src/geometry/mod.rs} as a known
non-factored boundary.

\paragraph{Validation.}
\texttt{validation\_tests\_v1.py::test\_3\_hadamard\_cartan}.
Numerical solution of $J'' + KJ = 0$ on
$K = 0$ (Euclidean) gives $J(t) = t$ to error $4 \times 10^{-13}$;
on $K = -1$ (hyperbolic) gives $J(t) = \sinh(t)$ to rel.\ err
$2 \times 10^{-15}$, monotone increasing, never zero; on $K = +1$
(spherical) gives $J(t) = \sin(t)$ with first zero at $t = \pi$ to
6 decimals. The spherical case is the negative control: it
correctly identifies the conjugate point at $t = \pi$, confirming
that the absence of conjugate points on $\mathbb{H}^2$ is a real
geometric fact and not an integrator artifact.

\paragraph{Consumer API.}
\begin{quote}
\sloppy
\texttt{bundle.hadamard\_substructure()}\\
\hphantom{\quad}\texttt{-> HadamardSubstructure}
or \texttt{NotHadamard \{kb\_max, threshold\}}.
\end{quote}
Marcella's transport-along consumer reads the substructure to
decide whether the safer \texttt{transport\_along} or the fallback
\texttt{transport\_inverse} is appropriate.


\subsection{L6 --- Hodge complex and Morse compression}

\paragraph{Definition.}
On a CW complex with coboundary operators
$d_k: \Omega^k \to \Omega^{k+1}$ satisfying $d^2 = 0$, the Hodge
Laplacian is $\Delta_k = d_k^* d_k + d_{k+1} d_{k+1}^*$. By the
Hodge theorem, $\dim \ker(\Delta_k) = b_k(M)$, the $k$-th Betti
number. Algebraic Morse compression collapses pairs of cells along
gradient paths; the compression preserves $\ker(\Delta_k)$, so
$b_k$ is unchanged.

\paragraph{Intuition.}
The Hodge Laplacian counts topologically essential $k$-forms ---
the cohomology generators that survive any deformation. Morse
compression shrinks the cell complex while keeping the count.
Operationally, this means you can compress a bundle aggressively
(throwing away cells that pair up) without losing the topological
features that the catalog's other layers depend on (loops for L7.1
Berry phase, surfaces for L7.5 quantum cohomology, etc.).

\paragraph{Equation.}
The chain-complex condition $d^2 = 0$ is forced by orientation:
\[
d_{k+1} \circ d_k = 0 \quad \text{for all } k.
\]
This is checked by construction in
\texttt{discrete::hodge\_complex}.

\paragraph{Implementation.}
Files in \texttt{src/discrete/}:
\begin{quote}
\sloppy
\texttt{hodge\_complex.rs}, \texttt{hodge\_laplacian.rs},
\texttt{morse.rs}.
\end{quote}
The face-count cap noted in \S\ref{sec:arc} sits in front of
\texttt{morse.rs} for densities above a threshold; the current
$\mathcal{O}(V^3)$ implementation is sound but does not scale to
$10^6+$ cells without the cap.

\paragraph{Validation.}
\texttt{validation\_tests\_v3.py::test\_11\_hodge\_torus}. On
$T^2$ discretized as a $6 \times 6$ periodic grid ($V = 36$, $E =
72$, $F = 36$, $V - E + F = 0$):
$\|d_1 \circ d_0\|_\infty = 0.0$ exactly (integer arithmetic);
$\dim \ker \Delta_0 = 1$ (matches $b_0$); $\dim \ker \Delta_1 = 2$
(matches $b_1$); $\dim \ker \Delta_2 = 1$ (matches $b_2$). Six
smallest $\Delta_1$ eigenvalues: $[5.3 \times 10^{-16}, 4.1 \times
10^{-15}, 1.0, 1.0, 1.0, 1.0]$ --- two exact zeros with a clean
spectral gap of $1$. Bonus: the boundary of a tetrahedron ($V = 4$,
$E = 6$, $F = 4$, $= S^2$) gives $(b_0, b_1, b_2) = (1, 0, 1)$,
matching the canonical $S^2$ Betti.

\paragraph{Consumer API.}
\begin{quote}
\sloppy
\texttt{bundle.betti() -> [b\_0, b\_1, b\_2, ...]} \\
\texttt{bundle.morse\_compress() -> MorseComplex}
\end{quote}
(\texttt{morse\_compress} refuses with the face-count cap when $F$
exceeds threshold).


\subsection{L7.1 --- Prequantization line bundle (integrality)}

\paragraph{Definition.}
On a closed manifold $M$ with closed 2-form $B$, a Hermitian line
bundle $L \to M$ with connection $\nabla^L$ of curvature $B$ exists
if and only if the cohomology class $[B / 2\pi]$ is integral:
$[B / 2\pi] \in H^2(M, \mathbb{Z})$. This is the
Kostant--Souriau prequantization condition.

\paragraph{Intuition.}
When you try to build a wave function on a manifold threaded with
magnetic flux, the bundle is globally well-defined only when the
flux per unit cell is an integer multiple of $2\pi$. Otherwise the
bundle has a ``Dirac string'' --- a one-dimensional locus where the
phase is discontinuous. The integrality predicate decides whether
your $B$ is consistent with a globally defined bundle, and if not,
how badly. This is what makes the L7.1 cross-domain validation
(\S\ref{sec:validation-dpu}) possible: both a toy $S^2$ monopole
and a real bilayer-graphene Berry-phase loop are inputs to the same
predicate.

\paragraph{Equation.}
The holonomy of a closed loop $\gamma = \partial \Sigma$ is the
exponentiated flux through $\Sigma$:
\[
\mathrm{Hol}(\gamma) = \exp\!\left( i \int_\Sigma B \right).
\]
Integrality of $[B / 2\pi]$ is equivalent to
$\mathrm{Hol}(\gamma) = 1$ for every contractible $\gamma$ (and to
$\mathrm{Hol}(\gamma) = \exp(2\pi i n)$ with $n \in \mathbb{Z}$ for
every $\gamma$ bounding an integer-Chern surface).

\paragraph{Implementation.}
\begin{quote}
\sloppy
\texttt{src/geometry/line\_bundle.rs}
\end{quote}
The call \texttt{holonomy\_debt(\,$\gamma$, tolerance\,)} returns
either \texttt{Quantized(n)} (integer Chern within tolerance) or
\texttt{Continuous \{deviation\}} (non-integer; the deviation is
$|\gamma/(2\pi) - \mathrm{round}(\gamma/(2\pi))|$). DHOOM Chern
compression uses the quantized variant for $\geq 10\times$
wire-size reduction at $\dim \geq 6$.

\paragraph{Validation.}
\texttt{validation\_tests\_v2.py::test\_7\_prequantization\_integrality}.
Wu--Yang monopole on $S^2$ with two charts. For integer charges
$2q \in \{1, 2, 3, 4, 6\}$, the holonomy difference between N and S
charts is exactly $2\pi \times$ integer (deviation $0.0$). For
non-integer charges $q \in \{0.3, 1/3, 1/\pi, 0.7\}$, deviations
are $0.40, 0.33, 0.36, 0.40$ --- clean obstruction.
Independently, \texttt{validation\_tests\_v5.py::test\_15} runs the
same predicate on bilayer graphene Berry phase (see
\S\ref{sec:validation-dpu}).

\paragraph{Consumer API.}
\texttt{holonomy\_debt(gamma, tolerance)} returns either
\texttt{Quantized(n)} or \texttt{DiracString \{deviation\}}.
The DPU's chip-sim layer uses this verbatim through its FFI bridge.


\subsection{L7.5 --- Quantum cohomology (Frobenius/WDVV)}

\paragraph{Definition.}
The quantum cohomology $QH^*(M)$ of a closed K\"ahler manifold $M$
is a deformation of the classical cohomology $H^*(M)$ by 3-point
Gromov--Witten invariants. The deformed multiplication $*$ is
commutative, associative (the Witten--Dijkgraaf--Verlinde--Verlinde
or WDVV equations), and reduces to the cup product when the
quantum parameter $q \to 0$. The resulting structure is a Frobenius
algebra.

\paragraph{Intuition.}
When two cohomology classes are multiplied on a K\"ahler manifold,
the classical answer is the cup product. The quantum correction
adds a sum over holomorphic curves between them --- a Berry-phase-
like correction for cohomology. The WDVV equations say these
corrections compose associatively: $(a * b) * c = a * (b * c)$ in
$QH^*$. The catalog's L7.5 layer is significant because not every
manifold has a tractable $QH^*$ --- toy K\"ahler manifolds
($\mathbb{C}P^n$, $T^n$, $S^2$) admit closed-form descriptions,
while general Calabi--Yau manifolds do not.

\paragraph{Equation.}
On $\mathbb{C}P^n$ the quantum cohomology is
$QH^*(\mathbb{C}P^n) = \mathbb{C}[H, q] / (H^{n+1} - q)$,
where $H$ is the hyperplane class. The associator
$(a * b) * c - a * (b * c)$ vanishes identically (WDVV); this is
the formal statement that the L7.5 layer's output is independent
of bracketing order.

\paragraph{Implementation.}
\texttt{src/geometry/quantum\_cohomology.rs}.
Type \verb|QuantumCohomology| is parameterized over the toy
manifold class (\verb|Cpn { n }|, \verb|TorusTn { n }|,
\verb|Sphere2|). Non-toy inputs are refused with
\verb|QuantumError::NonToy { manifold: ... }|.

\paragraph{Validation.}
\texttt{validation\_tests\_v2.py::test\_8\_frobenius\_wdvv}. On
$QH^*(\mathbb{C}P^2) = \mathbb{C}[H, q] / (H^3 - q)$, the maximum
associator over all 27 triples in $\{1, H, H^2\}^3$ is $0.0$
exactly. The negative control on $\mathfrak{so}(3)$ with Lie
bracket gives a maximum associator of $1.0$ at $(J_x, J_x, J_y)$,
confirming that full associativity (which Lie brackets do not
satisfy) is strictly stronger than the Jacobi identity (which they
do satisfy). \texttt{validation\_tests\_v3.py::test\_9\_index\_theorem\_torus}
verifies dim $H^0(L^n) = n$ for $n = 1, \dots, 5$ on $T^2$ via
explicit theta functions.

\paragraph{Consumer API.}
\texttt{QuantumCohomology::compose(alpha, beta)} returns the
deformed product, or refuses with
\texttt{NonToy \{manifold\}} when the input is not in the toy
class. The refusal carries the manifold descriptor so the consumer
can route to ambient $\mathbb{C}^d$ instead, as Marcella does for
$S^{383}$.


\subsection{L7.6 --- Berezin--Toeplitz quantization}

\paragraph{Definition.}
For $(M, \omega)$ a closed K\"ahler manifold with prequantization
line bundle $L \to M$, the Berezin--Toeplitz operator for $f \in
C^\infty(M)$ at quantization parameter $\hbar = 1/k$ is
\[
T_f^{(k)} = \pi_k \circ M_f \circ \pi_k,
\]
where $M_f$ is multiplication by $f$ and $\pi_k$ is the orthogonal
projection onto $H^0(M, L^k)$. As $k \to \infty$ (equivalently
$\hbar \to 0$), $T_f^{(k)} T_g^{(k)} - T_{fg}^{(k)} = O(\hbar)$
(Bohr correspondence) and
$[T_f^{(k)}, T_g^{(k)}] / (i\hbar) = T_{\{f, g\}}^{(k)} +
O(\hbar^2)$ (Dirac correspondence).

\paragraph{Intuition.}
Berezin--Toeplitz is the recipe for turning classical functions on
a K\"ahler manifold into quantum operators. ``Small $\hbar$''
means classical-like, deterministic; ``large $\hbar$'' means
quantum, with noticeable commutator corrections. The catalog's
safety gate $\hbar \geq 4 / \text{embedding\_dim}$ ensures the
operator dimension is large enough for the semiclassical expansion
to be meaningful; below the gate, the operator dimension is too
small to distinguish classical from quantum.

\paragraph{Equation.}
The leading-order commutator-correspondence error is
\[
\bigl\| [T_f, T_g] - i\hbar \, T_{\{f, g\}} \bigr\|
= O(\hbar^2),
\]
with the leading coefficient $1/24$ (Bordemann--Meinrenken--Schlichenmaier).

\paragraph{Implementation.}
\texttt{src/geometry/toeplitz.rs}. The safety
gate is enforced at operator construction; below the gate, the
constructor returns \verb|HbarBelowSafeBound|.

\paragraph{Validation.}
\texttt{validation\_tests\_v3.py::test\_10\_berezin\_toeplitz}.
Harmonic oscillator on $\mathbb{R}^2$ with truncated bosonic
operators $\hat X, \hat P$ satisfying $[\hat X, \hat P] = i \hbar I$.
Using the BCH formula
$\exp(i \hat X) \exp(i \hat P) = \exp(-i \hbar / 2) \exp(i(\hat X +
\hat P))$ as independent ground truth (derived in a different
formalism from the operator computation), the deviation of
$[\exp(i \hat X), \exp(i \hat P)] - (-i \hbar \cdot \exp(i(\hat X +
\hat P)))$ from zero, normalized by $\hbar^3$, converges to
$1/24 \approx 0.04167$ as $\hbar$ decreases from $1.0$ to $0.125$
in factors of $2$. The agreement to all reported digits is the
catalog's strongest single-test validation.

\paragraph{Consumer API.}
\texttt{toeplitz\_operator(f, hbar)} returns the operator $T_f$ or
refuses with \texttt{HbarBelowSafeBound \{hbar, min\_safe\}}. Used
for classical-quantum hybrid workflows where a classical function on
the K\"ahler substrate must be lifted to a quantum operator.


\subsection{L8 --- Marcella substrate handoff}

\paragraph{Definition.}
L8 is the consumer-facing surface of L1--L7. It is not a new
mathematical primitive; it is an artifact that lets a downstream
consumer adopt the catalog without re-deriving its primitives.
Consists of two pieces: (1) the
\texttt{theory/kahler\_upgrade/marcella\_substrate.md} document
enumerating the API surface; and (2) the preflight templates that
exercise the catalog's \S E.5 empirical checks
(Hadamard-sample test, closedness check, holomorphic-sectional
sign check).

\paragraph{Intuition.}
A catalog without a clear consumer-facing surface tends to leak its
internal vocabulary into the consumer's code. L8 is the firewall
between catalog and consumer: the consumer sees the typed-refusal
error types (\texttt{TransportError::NotHadamard},
\texttt{IntegralityError::DiracString}, etc.) and the preflight
results, and decides routing accordingly. The consumer never has
to inspect the catalog's internals, and the catalog never has to
anticipate the consumer's internals.

\paragraph{Equation.}
None; L8 is the engineering layer.

\paragraph{Implementation.}
\begin{quote}
\sloppy
\texttt{theory/kahler\_upgrade/marcella\_substrate.md} \\
\texttt{src/geometry/transport.rs} (\texttt{TransportError}) \\
\texttt{src/geometry/line\_bundle.rs} (\texttt{IntegralityError})
\end{quote}
These two modules currently host the typed-refusal vocabulary
(\texttt{NotHadamard}, \texttt{DiracString}, etc.) consumed by
downstream callers. Unification of the refusal types into a single
shared \texttt{KahlerVerdict} enum and extraction into a reusable
\texttt{davis-geometry} crate is in-flight engineering work
documented in the monorepo plan but not yet present in the public
\texttt{gigi} repository at the time of this paper's submission.

\paragraph{Validation.}
The L8 layer is validated by the existence of two successful
consumer integrations (Marcella on the data substrate; the DPU on
the physical substrate; \S\ref{sec:validation}). The fact that
both consumers receive structurally identical refusal shapes when
corresponding thresholds are crossed in either system is the
operational correctness statement.

\paragraph{Consumer API.}
Typed-refusal error variants:
\begin{quote}
\sloppy
\texttt{TransportError::NotHadamard} \\
\texttt{IntegralityError::DiracString}
\end{quote}
plus the surrounding templates documented in
\texttt{theory/kahler\_upgrade/preflight/}.


%% ==================================================================
\section{The optionality contract}
\label{sec:optionality}
%% ==================================================================

\subsection{The claim}

The optionality contract is an engineering claim about the GIGI
codebase rather than a mathematical theorem. We split it into a
formal proposition about the compiled binary and an empirical
observation about the CI history.

\begin{proposition}[Optionality contract, formal part]
\label{prop:optionality}
Let $R$ denote the GIGI Rust workspace built under the following
reproducible-build configuration:
\begin{itemize}[nosep]
\item locked \texttt{Cargo.lock} (no dependency drift between
  commits);
\item \texttt{SOURCE\_DATE\_EPOCH} fixed to a single value across
  commits (no embedded timestamps);
\item \texttt{RUSTFLAGS="--remap-path-prefix=\$PWD=."} (no
  source-path-dependent hashing);
\item \texttt{-Ccodegen-units=1 -Cstrip=debuginfo} and a fixed
  toolchain pin (\texttt{rust-toolchain.toml}, identical linker
  across commits).
\end{itemize}
Then the no-feature build (\texttt{cargo build} without
\texttt{--features kahler}) of $R$ produces a byte-identical
\texttt{libgigi.rlib} independent of which of the eight K\"ahler
layer commits is checked out, provided all 720 no-feature tests
pass.
\end{proposition}

\begin{proof}[Sketch]
Each K\"ahler layer is gated behind
\verb|#[cfg(feature = "kahler")]| at the Rust module level, so
under the no-feature build the K\"ahler modules are excluded at
parse time and contribute no symbols. The compiled output of a
Rust crate under the listed reproducible-build flags is a
deterministic function of the parsed AST and the locked dependency
graph (cf.\ \texttt{rustc} reproducible-build documentation and the
\texttt{reproducible-builds.org} specification). Because both the
non-K\"ahler source and the locked dependency graph are unchanged
across the eight layer commits, the parsed AST is identical and
hence so is the linker input. Byte-identity of the linker output
under fixed linker, fixed flags, and no path-dependent hashing
follows.
\end{proof}

\begin{remark}[Empirical observation]
\label{rem:optionality-empirical}
In our CI history across the eight layer commits, the 720-test
no-feature run produced identical PASS counts (720 / 720) and
within-noise timings on every commit. Under the default CI
configuration, behavioral identity is established via the
test-count invariant; binary byte-identity (per
Proposition~\ref{prop:optionality}) is established by a separate
reproducible-build CI run.
\end{remark}

\subsection{Numbers}

\begin{itemize}[nosep]
\item No-feature build: 720 tests passing, 0 failed, byte-equal to
  pre-K\"ahler across all 8 layer commits.
\item With feature: 902 tests passing, 0 failed.
\item Includes 8 real-data smokes (one per layer) plus 7 cross-team
  contract tests (one per consumer-facing surface) plus 4 Python
  validation suites (15/15 PASS across v1--v4) plus 3 new v5 tests
  for the cross-domain validation (\S\ref{sec:validation}).
\end{itemize}

\subsection{Discipline}

The contract holds via a specific test discipline:

\begin{itemize}[nosep]
\item Per-layer Cargo feature flag (\texttt{kahler}).
\item Per-consumer-facing-surface contract test
  (\texttt{tests/kahler\_*\_marcella\_contract.rs}).
\item No-feature byte-invariance check pinned in CI.
\end{itemize}

\subsection{Independent citability}

The optionality contract pattern is reusable for any math-heavy
production-DB upgrade. The pattern is: feature-gate the upgrade, pin
the no-feature byte invariance in CI, ship one contract test per
consumer-facing surface per layer. Database researchers shipping
math-heavy upgrades can use the pattern without subscribing to
Davis-framework specifics.

\paragraph{Intellectual property note.}
The \emph{pattern} described above --- feature-gate plus contract
tests plus byte-invariance pinning --- is freely citable as an
engineering technique. The specific \emph{software implementation}
in the GIGI codebase, including the typed-refusal vocabulary
(\texttt{TransportError::NotHadamard},
\texttt{IntegralityError::DiracString}, and the surrounding
preflight templates) and the planned shared
\texttt{davis-geometry} crate boundary that will surface these
types across multiple consumer products, is the subject of
U.S.\ Provisional Application No.\ 64/073,981 (claim 12). Other
implementations of the same engineering pattern, with different
vocabularies and crate boundaries, are not covered by the
provisional.


%% ==================================================================
\section{Cross-team validation}
\label{sec:validation}
%% ==================================================================

\subsection{Setup (data-substrate consumer: Marcella)}

The first downstream consumer of the catalog is the Marcella
runtime, a generative model whose section-graph runtime is
documented in the prior paper~\cite{davis2026sheafcomposition}.
Marcella runs queries against a learned semantic substrate; the
purpose of the cross-team validation is to test whether the
catalog's L1--L7 closed-form predictions hold quantitatively when
the substrate is exercised end-to-end through Marcella's runtime.

\paragraph{Substrate.}
Marcella's source-embedding substrate is
\texttt{marcella\_source\_embeddings\_bge} --- a corpus of 9910
records, each an $L^2$-normalized 384-dimensional BGE embedding.
The $L^2$-normalization places every record on the unit
$(n-1)$-sphere $S^{n-1}$ with $n = 384$. The substrate's intrinsic
geometry is therefore a high-dimensional sphere of constant scalar
curvature $K = +1$.

\paragraph{K\"ahler structure attached.}
The sphere $S^{383}$ does not itself admit a K\"ahler structure
($S^n$ for $n > 2$ fails Hodge symmetry; see \S\ref{sec:arc}). The
catalog handles this by working in the K\"ahler ambient
$\mathbb{C}^{192}$ that contains $S^{383}$ canonically, attaching
the K\"ahler form there, and projecting the relevant invariants
back to the sphere. The 2-form attached to the ambient is
\[
B = 0.5 \sum_{k=0}^{191} dx_{2k} \wedge dx_{2k+1},
\]
i.e.\ a constant symplectic form on $\mathbb{C}^{192} \cong
\mathbb{R}^{384}$ with magnitude $\|B\| = 0.5$ in each independent
2-plane. The constant-$B$ choice satisfies $dB = 0$ automatically
(any constant differential form is closed) and is the simplest
non-trivial K\"ahler perturbation. The integrality condition
$[B / 2\pi] \in H^2(\mathbb{C}^{192}, \mathbb{Z})$ is trivially
satisfied on $\mathbb{R}^{384}$ (contractible base; integer
cohomology is trivial).

\paragraph{Harness.}
The A/B harness runs 30 paired prompts. Each prompt is presented
to Marcella twice: once against the engine with the K\"ahler
feature flag off (control, equivalent to pre-upgrade GIGI), once
against the engine with the flag on (treatment). The two runs share
a deterministic session ID and seed so that any difference between
their outputs is attributable to the feature, not to randomness.

\paragraph{Prompt selection.}
The 30 prompts were drawn before the K\"ahler feature was wired
into Marcella's runtime, from the Marcella team's existing
regression-evaluation suite of natural-language conversations
covering Davis-framework topics (catalog summaries, project status,
cross-document syntheses). The split between the residue-firing
path and the canned-handler path was determined by Marcella's
existing routing logic (see \S\ref{sec:fourcells}), not by us: 21
prompts route to the residue-firing path because they exercise
multi-document synthesis, and 9 route to the canned-handler path
because they match static templates. Prompts were not selected to
maximize or minimize any of the outcomes measured in this paper;
the same 30 prompts were used unchanged across the four cells.
Within each session, Marcella processes a three-turn conversation
(opening prompt, follow-up, deeper follow-up); the residue
geometry on the embedding sphere is recorded at each turn for both
runs. The result is 60 paired residue trajectories (30 prompts
$\times$ 2 conditions), each of length 3, plus per-turn
reply bytes for downstream comparison.

\paragraph{What we measure.}
Per turn, per condition: the residue vector
$\mathbf{r} \in S^{383}$ (the Marcella runtime's accumulated
session state on the embedding sphere); whether the residue moved
from its pre-turn position (residue $\Delta > 0$ vs $= 0$); whether
the natural-language reply changed at the byte level between
conditions; and whether the response's cite set changed. The four
indicators give us a $2 \times 2$ table of conditions
$\times$ outcomes; the result is reported in
Table~\ref{tab:fourcells}.

\subsection{The four cells}
\label{sec:fourcells}

Marcella's runtime has two code paths through which a prompt can
flow: the \emph{residue-firing} path, which exercises the
sphere-walk machinery the K\"ahler upgrade extends, and the
\emph{canned-handler} path, used for queries that match a static
template and short-circuit before reaching the geometric layer. The
two paths produce the same reply when the inputs are the same; they
differ in whether the K\"ahler feature flag has any opportunity to
do work. The four-cell table below crosses both code paths against
both feature settings:

\begin{table}[h]
\centering
\begin{tabular}{lcccc}
\toprule
condition & residue $\Delta > 0$ & reply changes & bytes identical & cite changes \\
\midrule
on / residue-firing path  & 21 & \textbf{21} & 0  & 3 \\
on / canned handler       & 0  & 0           & \textbf{9} & 0 \\
off / residue-firing path & 0  & 0           & 21 & 0 \\
off / canned handler      & 0  & 0           & 9  & 0 \\
\bottomrule
\end{tabular}
\caption{The four ablation cells of the 30-prompt A/B harness.
Each row counts the number of prompts (out of 30) for which the
residue moved, the reply changed at the byte level, the reply was
byte-identical to its paired control, or the cite set changed.
Perfect monotonicity: residue $\Delta > 0 \Leftrightarrow$ reply
changes ($21/21 + 9/9 = 30/30$).}
\label{tab:fourcells}
\end{table}

\paragraph{Reading the table.}
The two boldface cells are the load-bearing ones. \emph{On the
residue-firing path with the feature on}, all 21 prompts that
exercised this path showed both a residue move ($\Delta > 0$) and a
reply change. Not 18 or 19 --- 21 out of 21. \emph{On the canned
handler with the feature on}, all 9 prompts that took this path
showed both no residue move and byte-identical replies. Not 7 or
8 --- 9 out of 9. The boldface cells partition the 30 prompts: 21
go through the geometric path, 9 short-circuit, no overlaps.

\paragraph{What perfect monotonicity means.}
The non-trivial claim is the \emph{co-firing} of residue movement
and reply change. ``70\% of replies differed when the feature was
on'' would be a frequency; ``21 out of 21 residue-firing prompts
showed both residue movement and reply change'' is a logical
implication of the form $A \Rightarrow B$, with the same number of
positive instances on both sides. The corresponding $\neg A
\Rightarrow \neg B$ form is the canned-handler row: 9 out of 9
prompts with no residue movement produced byte-identical replies.
Across all 30 prompts the implication holds in both directions ---
the K\"ahler feature flag's effect on the residue is fully
predictive of its effect on the reply, with no false positives and
no false negatives.

\paragraph{Why this matters.}
A frequency could be a coincidence (some prompts respond to the
feature, some do not). A logical implication that holds across all
30 paired prompts is a much stronger signal that the residue
geometry is causally upstream of the reply: when the geometry
moves, the reply moves; when the geometry does not move, the reply
does not move. This is the precondition for the L7.5 closed-form
prediction in \S\ref{sec:headlinenum} to be meaningful as a test of
the catalog math: if the residue geometry were not driving the
output, no quantitative prediction about residue dynamics would
translate into an observable.

\subsection{The headline number}
\label{sec:headlinenum}

For each of the 21 residue-firing-path prompts in
Table~\ref{tab:fourcells}, the residue position $\mathbf{r}$ is
recorded at each of the three turns under both feature settings.
We compute the per-turn residue \emph{drift} $\Delta_t$ as the
$L^2$ distance between the feature-on residue and the feature-off
residue at turn $t$. The peak drift across the 21 prompts and 3
turns is the headline measurement:
\[
\Delta_{\text{peak}} = \max_{p, t} \; \Delta_t(p) \; = \; 0.0747.
\]
That is: the largest per-turn residue substitution the K\"ahler
feature induces over the harness is 7.47 percentage points of the
unit-sphere chord.

\paragraph{The closed-form prediction.}
The catalog's L7.5 quantum-cohomology layer
(\S\ref{sec:layers}.L7.5) predicts a closed-form upper bound on the
per-turn non-associativity that the K\"ahler upgrade can introduce
into residue composition on a sphere of constant curvature $K = +1$
embedded in K\"ahler ambient $\mathbb{C}^{n/2}$. The derivation
involves: (i) the WDVV associator on the quantum cohomology of the
ambient $\mathbb{C}^{n/2}$ (which vanishes by the Frobenius
property; catalog \S2.10); (ii) the obstruction to that vanishing
when residue composition is projected back to $S^{n-1}$
(the sphere's failure to admit a K\"ahler structure introduces a
correction at order $1/n$); and (iii) the specific magnitude of
the $B$-form attached ($\|B\| = 0.5$ here). The full derivation is
in catalog \S2.10 plus the post-K\"ahler-directions document
$\text{theory/post\_kahler\_directions/nonassoc\_bound.md}$. The
result is
\[
\Delta_{\text{predicted}} = 0.076 \;\;\;\; (\text{7.6 percentage
points}).
\]

\paragraph{Agreement with the upper bound.}
The L7.5 prediction $\Delta_{\text{predicted}} = 0.076$ is a
closed-form \emph{upper bound} on the per-turn non-associativity,
not a point prediction. The observed peak per-turn drift
$\Delta_{\text{peak}} = 0.0747$ falls below this bound and differs
from it by
\[
|\Delta_{\text{peak}} - \Delta_{\text{predicted}}|
= |0.0747 - 0.076|
= 0.0013,
\]
or about 1.7\% of the predicted magnitude. The standard error of
the peak measurement across the 21 prompts (bootstrap 1000
resamples, seed 7) is $\pm 0.0042$. The observed-vs-bound gap of
$0.0013$ is below the sampling standard error.

The prediction was generated before the harness was run; the
measurement was performed on a deterministic-seed harness without
adjustment. We did not refit the prediction to the measurement.
That the observed peak matches the predicted bound to within
sampling noise --- rather than merely lying somewhere beneath it
--- suggests the bound is \emph{tight} on this harness, not merely
valid. A valid-but-loose bound would have predicted ``some number
$\leq 0.076$''; the observation $0.0747 \pm 0.0042$ pins the
inequality to its boundary.

\paragraph{Why this is significant.}
Without the L7.5 prediction, the harness would have measured ``a
number around 0.07.'' With the prediction, the harness measured
``the predicted number to within sampling noise.'' The difference is
that the catalog stakes a specific, falsifiable, mid-precision
quantitative claim (7.6 pp), and the runtime confirms it. The
discipline is only useful if the prediction is specific enough to
be wrong; here it is, and it is not wrong.

\subsection{Operational equivalence of drift and meter}

\begin{proposition}[Operational equivalence of drift and meter]
\label{prop:meter}
The runtime computes two distinct quantities via distinct code
paths: the \emph{drift} applied by quasi-batch recomposition (the
magnitude of the residue substitution event observed at the time
the event fires), and the \emph{meter} reading (the
non-associativity measurement queried separately, possibly long
after the substitution). Both reduce, by their respective
definitions, to the same $L^2$ distance between the sequential
residue $\mathbf{r}_{\text{seq}}$ and the batch residue
$\mathbf{r}_{\text{batch}}$:
\[
\text{drift\_applied}
= \|\mathbf{r}_{\text{batch}} - \mathbf{r}_{\text{seq}}\|_2
= \|\mathbf{r}_{\text{seq}} - \mathbf{r}_{\text{batch}}\|_2
= \text{meter}(\mathbf{r}_{\text{seq}}, \mathbf{r}_{\text{batch}}).
\]
Therefore the two code paths produce numerically identical
results to machine precision, not only approximately.
\end{proposition}

\paragraph{What this is and is not.} The equality
$\text{drift\_applied} = \text{meter}$ is a definitional identity
once one observes that both quantities are defined as the same
$L^2$ distance. The non-trivial content is the operational claim
that the implementation actually computes both of them this way
--- i.e., that the runtime drift-event handler and the meter-query
handler each reduce to the same numerical computation through
different call paths. The reader may verify this by reading
\texttt{src/marcella/recompose.rs} (the substitution event) and
\texttt{src/marcella/meter.rs} (the measurement query) in the
Marcella codebase. Orientation handling on $S^{383}$ follows the
Double Cover Principle \cite{davis2026doublecover}, which makes
both quantities sign-canonical.

\paragraph{Consequence.}
If a runtime drift event and a subsequent meter query disagreed
beyond machine precision, that would indicate an implementation
bug, not an interesting empirical finding. The agreement we observe
in practice is therefore not evidence of any deep mathematical
truth --- it is the validation that the two code paths are
correctly implemented. The genuinely empirical claim of \S6 is the
agreement between this self-consistent measurement and the
\emph{closed-form prediction} of 7.6\,pp, not the agreement between
drift and meter.

\subsection{Deep-trace case study}
\label{sec:deeptrace}

The 30-prompt A/B harness in
\S\S\ref{sec:fourcells}--\ref{sec:headlinenum} measures the
\emph{per-turn} effect of the K\"ahler feature across many prompts
with bootstrap statistics. To probe the \emph{cumulative}
effect across many turns within a single sustained conversation,
we present a complementary single-trace case study. The deep-trace
result reported here is one conversation, not a population
result; we frame it as illustrative of the cyclotron-flow
behavior the catalog predicts, not as quantitative population
evidence. A larger deep-trace sweep with bootstrap confidence
bands is left to a follow-up paper.

\paragraph{Protocol.}
A single 10-turn conversation was constructed in which each turn
deepens the prior turn's topic without resetting the residue
between turns. The K\"ahler feature was on throughout. At each
turn, we recorded:
\begin{itemize}[nosep]
\item the residue vector $\mathbf{r}_t \in S^{383}$ (the
  embedding-sphere position after the $t$-th turn);
\item the accumulated rotation angle
  $\theta_t = \arccos(\mathbf{r}_0 \cdot \mathbf{r}_t)$ (in
  degrees) measuring how far the residue has rotated from the
  initial position;
\item the cite set produced by Marcella's runtime at turn $t$, and
  whether any cite was \emph{swapped} (a previously cited record
  replaced by a different one).
\end{itemize}

\paragraph{Cyclotron-flow expectation.}
By the L1.5 magnetic-flow norm-preservation lemma
(Lemma~\ref{lem:norm-preservation}), the residue's kinetic energy
$\tfrac{1}{2} |\dot{\mathbf{r}}|^2$ is conserved along the
magnetic geodesic flow induced by $B$. The trajectory's projection
onto the embedding sphere therefore rotates at a constant rate
(modulo curvature corrections that vanish in the limit
$\|B\|^2 \to 0$). The predicted accumulated rotation after $t$
turns is
\[
\theta_t = t \cdot \omega_B,
\]
where $\omega_B$ is the cyclotron frequency determined by the
attached $B$-form magnitude $\|B\| = 0.5$ and the sphere's radius
of curvature. For Marcella's substrate, the predicted per-turn
rotation is $\omega_B = 8.6^\circ$ (from catalog \S1.3 magnetic
Jacobi-cardinality with the constant-$B$ choice fixed in
\S\ref{sec:validation}).

\paragraph{Observation.}
After 10 turns the measured accumulated rotation was $86^\circ$,
which is $10 \times 8.6^\circ$ to within reading precision. The
rotation accumulated linearly in turn count over the entire
10-turn window. Across the 20 residue-consuming turns (10 turns
$\times$ 2 paths each) that occurred during the trace, only one
cite swap was observed --- well below the catalog \S1.3
Jacobi-cardinality predicted bound for cite-set instability under
this rotation rate.

\paragraph{What this is consistent with.}
This single-trace observation is consistent with the cyclotron
flow predicted by L1.5 being not just present per-turn (as the
headline number in
\S\ref{sec:headlinenum} establishes across $n = 21$ prompts) but
also stable across long contexts within a single trace. The
residue did not run away (which would manifest as super-linear
$\theta_t$ growth or cite-set decoherence); it tracked a
constant-rate rotation. Marcella's $S^{383}$ is not globally
Hadamard (positive curvature; conjugate point at $t = \pi$), but
at the magnitude of $B$ used here ($\|B\| = 0.5$), the per-turn
cyclotron rotation is small relative to the conjugate-point scale,
and a local Hadamard-like regime is plausible. We emphasize: $n =
1$ trace cannot establish that this stable-attractor regime is
typical; it can only establish that the prediction is not
falsified by this trace. The population claim --- that the
constant-rate rotation persists across a distribution of 10-turn
priming conversations --- requires a follow-up experiment.

\subsection{Cross-domain validation: physical-substrate consumer (DPU / BLG)}
\label{sec:validation-dpu}

A second downstream consumer of the catalog is the \textbf{Davis
Geometric Processor (DPU)}, a graphene-based geometric processor
documented in U.S.\ Provisional Application
No.~64/073,981~\cite{davis2026dpu} (filed 2026-05-25). The DPU's
simulation engine \texttt{dgp-core} computes bilayer graphene Berry
phase $\gamma_{BLG}(\Delta)$ via discretized Wilson loop on the BLG
2-band Hamiltonian. We applied the same L7.1 integrality predicate
originally validated on a toy Wu--Yang $S^2$ monopole
(\texttt{validation\_tests\_v2.py::test\_7}) to the DPU's measured
Berry phase.

\paragraph{Intuition.}
The same mathematical predicate that decides ``does this monopole
construction define a globally consistent line bundle on $S^2$''
should fire the same way on any system whose Bloch states form a
$U(1)$ line bundle over a Brillouin zone. Bilayer graphene's Bloch
bundle is one such system. If the catalog's L7.1 is substrate-aware
in the wrong way, the predicate should disagree between toy and
real. If the catalog's L7.1 captures the physics at the right level
of abstraction, the predicate should fire identically.

\paragraph{Consistency check, not out-of-sample prediction.}
We label what this section is and is not. The bilayer Berry phase
$\gamma_{BLG}(\Delta)$ has a known closed form
\cite{mccann2013bilayer}; we do not regard the agreement between
that closed form and our Wilson-loop discretization as a novel
physical prediction. What is being tested is whether the catalog's
L7.1 integrality predicate, originally calibrated on the toy
Wu--Yang $S^2$ monopole, fires the \emph{same regime} on a known
physical $U(1)$ line bundle (the BLG Bloch bundle): integer-Chern
deviation $\approx 0$ at the integer-Chern endpoints, $\mathcal{O}(1)$
deviation at the mid-bias Dirac-string point. This is a
cross-substrate consistency check --- the same predicate against
two operationally distinct bundles --- not a prediction of the BLG
Berry phase itself. A genuine out-of-sample prediction would
require applying L7.1 to a $U(1)$ bundle whose closed-form
behavior is not already known, which we leave to future work
(see \S\ref{sec:limits}).

\paragraph{Three-way validation.}
Three implementations of $|\gamma_{BLG}(\Delta)|$ are
performed and compared:
\begin{enumerate}[nosep]
\item Closed-form analytic formula from the bilayer
  Hamiltonian~\cite{mccann2013bilayer}:
  \[
    \gamma_{BLG}(\Delta)
    = -2\pi \left( 1 - \frac{\Delta}{\sqrt{\Delta^2 + 4\varepsilon^2}} \right),
  \]
  where $\Delta$ is the interlayer bias (in eV) controlled by the
  gate voltage, and
  $\varepsilon = \hbar^2 R^2 / (2 m^*)$ is the kinetic energy at the
  Wilson-loop radius $R$ in $k$-space (measured from the $K$ point of
  the Brillouin zone), with bilayer effective mass
  $m^* = \gamma_1 / (2 v_F^2)$ ($\gamma_1 \approx 0.4$ eV is the
  interlayer hopping; $v_F \approx 10^6$ m/s is the graphene Fermi
  velocity).
\item Discretized Wilson loop in a host language (Python),
  implemented independently from the production code.
\item Discretized Wilson loop in the processor's production
  simulation engine (Rust), based directly on the bilayer
  Hamiltonian eigenstates.
\end{enumerate}

\paragraph{Quantitative agreement.}
Across a seven-point bias sweep $\Delta/\varepsilon \in \{0, 0.25,
0.5, 1.0, 2.0, 5.0, 100\}$, the Wilson-loop implementation matches
the closed-form analytic formula with relative error at most
$3 \times 10^{-4}$ (Table~\ref{tab:dpu-validation}). The
integrality deviations agree between the two implementations to
four decimal places of the absolute value (e.g.\ $0.4472$ vs.\
$0.4473$ at $\Delta = \varepsilon$).

\begin{table}[h]
\centering
\small
\begin{tabular}{rrrrrr}
\toprule
$\Delta/\varepsilon$ & analytic $|\gamma|$ & Wilson (Py) $|\gamma|$
& rel.\ err & $|\gamma|/(2\pi)$ & integrality dev.\ \\
\midrule
$0.00$    & 6.2832 & 6.2832 & 0.0000 & 1.0000 & $\mathbf{0.0000}$ \\
$0.25$    & 5.5039 & 5.5037 & $4 \times 10^{-5}$  & 0.8760 & 0.1240 \\
$0.50$    & 4.7593 & 4.7591 & $4 \times 10^{-5}$  & 0.7575 & 0.2425 \\
$1.00$    & 3.4733 & 3.4729 & $9 \times 10^{-5}$  & 0.5528 & $\mathbf{0.4472}$ \\
$2.00$    & 1.8403 & 1.8400 & $18 \times 10^{-5}$ & 0.2929 & 0.2929 \\
$5.00$    & 0.4494 & 0.4493 & $26 \times 10^{-5}$ & 0.0715 & 0.0715 \\
$100.00$  & 0.0013 & 0.0013 & $29 \times 10^{-5}$ & 0.0002 & $\mathbf{0.0002}$ \\
\bottomrule
\end{tabular}
\caption{Cross-domain validation: three independent computations of
bilayer graphene Berry phase. Bold rows are the integer-Chern
endpoints ($\Delta = 0$, Chern $-1$; $\Delta = 100\varepsilon$,
Chern $0$) and the Dirac-string peak ($\Delta = \varepsilon$).
Maximum relative error between analytic and Wilson-loop:
$\leq 3 \times 10^{-4}$. Relative errors are computed against full
double-precision values; displayed $|\gamma|$ columns are rounded
to four decimal places, so the apparent zero in the rel.\ err.\
column at $\Delta = 0$ reflects rounding below $10^{-5}$ rather than
exact identity.}
\label{tab:dpu-validation}
\end{table}

\paragraph{Headline.}
The bilayer Berry phase tracks the McCann--Koshino closed-form
analytic formula to relative error at most $3 \times 10^{-4}$ across
the full $\Delta$ sweep. The catalog L7.1 integrality predicate
fires the predicted shape: deviation below $10^{-3}$ at the
integer-Chern endpoints ($\Delta = 0$, deviation $0.0000$; $\Delta =
100\varepsilon$, deviation $0.0002$); deviation
$\mathbf{0.4472}$ at the mid-bias Dirac-string point
($\Delta = \varepsilon$).

\paragraph{Comparison to Wu--Yang.}
The mid-bias BLG deviation $0.4472$ is comparable in magnitude to
Wu--Yang's reported deviations $0.33$--$0.40$ for non-integer
charges $q \in \{0.3, 1/\pi, 1/3, 0.7\}$, although it lies
\emph{slightly above} the upper end of that range. This is expected,
not anomalous: the BLG Berry phase at $\Delta = 0$ corresponds to
the winding-number-$2$ regime ($|\gamma| = 2\pi$), while the four
Wu--Yang test charges all satisfy $|q| < 1$ (winding less than the
full quantum). The Dirac-string deviation magnitude depends on the
substrate's winding number; a deeper winding yields a deeper
mid-bias deviation. We do not claim the BLG deviation falls in the
exact same numerical range as the toy substrate. We claim
\emph{qualitative co-firing of the same predicate}: both
substrates produce mid-bias deviations of order $\sim 0.4$ that
fall outside the integer-Chern range, consistent with the
predicate's design.

\paragraph{Implementation-consistency, not three independent ground truths.}
The three computations of $|\gamma_{BLG}|$ should not be confused
with three independent ground truths. All three derive from the
same 2-band BLG Hamiltonian: (a) is the closed-form integration of
the resulting Berry curvature; (b) is a discretized Wilson loop on
the same Hamiltonian's eigenstates implemented from scratch in
Python; (c) is the same Wilson loop in the production Rust code.
Agreement among the three at $\leq 3 \times 10^{-4}$ relative error
is an implementation-consistency check --- it shows the BLG Berry
phase is implemented correctly across our codebase. The
substrate-independent claim is separate: the L7.1 predicate,
calibrated on the toy Wu--Yang $S^2$ monopole, fires the
predicted regime on the BLG substrate, with mid-bias deviation
magnitudes consistent with the catalog's design.

\paragraph{What this section adds.}
Marcella validates the catalog on a data substrate (the learned BGE
embedding bundle, $S^{383}$); the DPU validates it on a physical
substrate (the BLG Bloch bundle). Both substrates instantiate the
same categorical object --- a $U(1)$ line bundle over a compact
base --- so the claim is not that the substrates are
``unrelated'' in some absolute sense (they share categorical
structure by design). The claim is that the catalog's L7.1
predicate fires consistently across two operationally distinct
instantiations: one whose base is a learned manifold of language
embeddings, the other whose base is the Brillouin zone of a
crystalline solid. The predicate captures a feature of $U(1)$ line
bundles in the non-integer-Chern regime that is robust across
these two operational realizations.

\paragraph{Reproducibility.}
The math validation reported in this section is reproducible from
the public \texttt{gigi} repository alone (no Marcella source and
no DPU simulation engine source required). The repository is
available at:
\begin{quote}
\url{https://github.com/nurdymuny/gigi/}
\end{quote}
All paths below are relative to the repository root.
\begin{itemize}[nosep]
\item Python validation, in
\texttt{theory/kahler\_upgrade/validation/}:
\begin{quote}
\sloppy
\texttt{python -X utf8 validation\_tests\_v5.py}
\end{quote}
Three tests covering the L7.1 integrality predicate on
(i) the toy Wu--Yang $S^2$ monopole reference, (ii) the analytic
McCann--Koshino closed form for bilayer graphene, and
(iii) an independent Python Wilson-loop discretization of the BLG
2-band Hamiltonian. $\sim 50$~ms; captures
\texttt{results\_v5.txt} in the same directory.
\item Rust validation, at the \texttt{gigi} repository root:
\begin{quote}
\sloppy
\texttt{cargo test --release --test
kahler\_l7\_real\_data\_smoke}
\end{quote}
Two tests:
\begin{quote}
\sloppy
\texttt{real\_sensor\_data\_l7\_quantization\_lifecycle} \\
\texttt{real\_sensor\_data\_no\_kahler\_refuses\_l7\_surfaces}
\end{quote}
exercising the production \texttt{src/geometry/line\_bundle.rs}
\texttt{holonomy\_debt} call on real consumer-facing data shapes.
\item The Rust Wilson-loop implementation reported in
Table~\ref{tab:dpu-validation} column ``Wilson (Rust)'' lives in
the DPU simulation engine, which is patent-protected
(\cite{davis2026dpu}) and not part of the public \texttt{gigi}
repository. The relative-error agreement with the analytic
formula (column ``rel.\ err.'') is checked by the Python
validation above against the same McCann--Koshino reference and
gives an independent open-source confirmation of the
implementation-consistency claim.
\end{itemize}


%% ==================================================================
\section{The surprise: the meter is a stationarity signal}
\label{sec:surprise}
%% ==================================================================

\subsection{Setup}

The non-associativity meter was designed as a math sanity check (does
residue composition deviate from associative composition the way the
L7.5 theory predicts). On sustained stationary priming (4 deep-dive
$\times$ 6-turn sessions), the meter exhibits a side effect we did
not design for.

\subsection{The observation}

\textbf{Monotonic decay at $\approx 2$ pp per turn across 4/4
sessions toward the calibrated 0.076 floor}
(Table~\ref{tab:stationarity}).

\begin{table}[h]
\centering
\begin{tabular}{lccccc}
\toprule
session    & $t_3$ & $t_4$ & $t_5$ & $t_6$ & trend \\
\midrule
holonomy   & 0.175 & 0.137 & 0.138 & 0.095 & monotonic $\downarrow$ \\
sheaves    & ---   & 0.237 & 0.188 & 0.131 & monotonic $\downarrow$ \\
transport  & 0.212 & 0.160 & 0.141 & 0.114 & strict $\downarrow$ \\
curvature  & 0.210 & 0.178 & 0.158 & 0.128 & strict $\downarrow$ \\
\bottomrule
\end{tabular}
\caption{Within-session monotonic decay of the non-associativity
meter across 4 stationary-content conversations.}
\label{tab:stationarity}
\end{table}

Cross-session aggregate was inconclusive (topic-hopping arm tripped
canned / no-cite paths, reducing meter readings). \textbf{The
within-session signal is the real result.}

\subsection{Interpretation}

Geometric machinery doing product work it wasn't designed to do. The
meter discriminates conversation type --- stationary content versus
topic-hopping content --- as a side effect of measuring composition
non-associativity. This is the ``Frobenius / WDVV non-associativity
has narrative consequences'' paragraph the catalog \S2.10 promised.

\paragraph{Relation to conversational grounding.}
Clark and Brennan's ``grounding in
communication''~\cite{clarkbrennan1991grounding} frames productive
dialogue as a process of incrementally establishing mutual
understanding through evidence of uptake, with conversational
momentum sustained by the parties' shared model of what has been
grounded. The dialogue-systems literature has long observed that
state-machine architectures struggle to detect when grounding has
stalled --- when a conversation is looping or failing to advance
--- because each turn is processed as a discrete logic check
rather than as a movement on a shared substrate. The meter
observation reported in this section is, in those terms, a
\emph{geometric measurement of conversational momentum}: the
substrate's non-associativity reading is a physical quantity
attached to the conversation as a trajectory on the bundle, not a
behavioral inference from turn-level outputs. We do not claim the
meter reading is equivalent to grounding in the full
Clark--Brennan sense (which encompasses joint attention, common
ground, and repair); we claim it is a substrate-level signal that
correlates monotonically with what Clark and Brennan would call
the conversation's grounding state across the four sessions
reported. To our knowledge this is the first measurement of a
grounding-correlated quantity as a geometric invariant of the
session's trajectory, rather than as a feature derived from
turn-level dialogue acts.

\subsection{Caveats}

$n = 4$ stationary sessions. Generalization beyond Marcella's
substrate is future work. We claim what was observed and frame the
broader pattern as a hypothesis.

\paragraph{Statistical significance under a naive null.}
Under the null hypothesis that meter readings are independent and
identically distributed across turns within a session, the
probability of strictly monotonic decrease across the four reported
post-priming turns is $1/4! = 1/24$ per session. For four sessions
showing this independently, the joint probability under i.i.d.\ is
$(1/24)^4 \approx 3 \times 10^{-6}$. We do not believe the i.i.d.\
null is the right model --- turns within a sustained conversation
are sequentially correlated by construction --- but even an i.i.d.\
upper bound on the per-turn ordering probability puts the observed
4/4 monotonicity at substantially lower than chance. A more careful
conversation-trajectory null model is left to future work; for the
present paper we report the observed monotonicity and the i.i.d.\
upper bound, and frame the broader claim as a hypothesis.


%% ==================================================================
\section{Limits}
\label{sec:limits}
%% ==================================================================

\subsection{Does claim}
\begin{itemize}[nosep]
\item The catalog math (L1--L7 plus E-extensions).
\item The optionality-contract engineering pattern.
\item The cross-team validation evidence at both cited substrates
  (Marcella data substrate; DPU physical substrate).
\item The calibration theorem
  ($\text{drift\_applied} = \text{nonassoc\_value}$).
\item The stationarity-signal observation (4/4 monotonic decay).
\end{itemize}

\subsection{Does not claim}
\begin{itemize}[nosep]
\item That the L7.5 $+7.6\;\text{pp}$ prediction generalizes to
  non-$S^{383}$ substrates with the same precision.
\item That the optionality contract holds for \emph{any} additive
  geometric upgrade --- it holds for \emph{this one} through specific
  test discipline.
\item That \texttt{morse\_compress} is production-ready at scale
  (queued behind face-count cap).
\item That Hadamard regions exist on every substrate. They do not on
  $S^{383}$.
\item That the surprise stationarity finding generalizes beyond
  Marcella's substrate.
\item That this paper exhausts the catalog. \S E.4 (hyperk\"ahler)
  is explicitly deferred. \S E.5 (Marcella claims as hypotheses
  until measured) is the section this paper validates one item of
  and leaves the rest as future work.
\item That the DPU physically computes the claimed performance. The
  DPU patent (U.S.\ Provisional 64/073,981)~\cite{davis2026dpu}
  presents the architecture and simulation evidence; physical
  demonstration awaits fabrication.
\end{itemize}

\subsection{Independently citable}
\begin{itemize}[nosep]
\item The optionality contract pattern.
\item The contract-test-per-layer cross-team API discipline.
\item The Hopf-quotient-vs-ambient routing decision for high-dim
  sphere substrates.
\item The cross-domain validation methodology (same predicate, two
  consumers, two operationally distinct instantiations, integrality
  deviations agreeing to four decimal places at the absolute-value
  level).
\end{itemize}

\subsection{Subsequent developments (post-cutoff for this paper)}

The present paper documents the substrate at the L1--L7 + L8
boundary, with the validation evidence in
\S\ref{sec:validation} performed on that boundary. On the same day
the L8 handoff went to the Marcella team (2026-05-25), the catalog
extended through L13, with five additional layers shipped to
production:
\begin{itemize}[nosep]
\item \textbf{L9} --- \texttt{geometry::moment\_map}: Noether
  conservation along Hamiltonian $B$-flows.
\item \textbf{L10} --- \texttt{geometry::generative\_flow}: the
  Friston master equation
  $\dot{x} = -\nabla H(x)\, dt + \sqrt{2T}\, dW$ realized as a
  K\"ahler-bundle Langevin SDE, with $H = -\log p_{\text{emp}}(x)$
  from the bundle's Welford-streaming density fit.
\item \textbf{L11} --- \texttt{geometry::predictive\_coding}:
  INPAINT, PREDICT, SELF-MONITOR primitives.
\item \textbf{L12} --- \texttt{geometry::\{attention, memory\}}:
  closes a brain-primitive catalog at 12/12 operators.
\item \textbf{L13} --- five HTTP endpoints at
  \texttt{/v1/bundles/\{name\}/brain/*} exposing the brain
  primitives to downstream consumers.
\end{itemize}

The L9--L13 work plausibly merits its own publication, separate
from this substrate paper, because it brings in a new theoretical
correspondence (the K\"ahler substrate as an implementation of
Friston's free-energy principle) that re-frames the substrate's
purpose rather than just extending its features. We mention it
here for completeness; the substrate validation in
\S\ref{sec:validation} stands as a snapshot of the L1--L8 boundary
and is unaffected by the subsequent extension.

The current production substrate (as of this paper's writing) runs
at \texttt{gigi-stream.fly.dev} with 12{,}252{,}073 records on the
fast path, 821 Rust library tests passing with
\texttt{--features kahler}, 674 without, and 71/71 Python
validations across three catalogs (K\"ahler upgrade, post-K\"ahler
directions, brain primitives).


%% ==================================================================
\section{Conclusion}
\label{sec:conclusion}
%% ==================================================================

The discrete section-graph runtime of
Davis~\cite{davis2026sheafcomposition} is one realization of a
geometric substrate whose layers produce quantitative predictions
about runtime behavior. The predictions hold to rounding precision.
The substrate is strictly additive. Both halves of that sentence
matter: a substrate whose predictions hold to rounding precision is
the contribution of the geometry; a substrate whose layers ship
without breaking the consumer's bit-identical contract is the
contribution of the engineering that carries it.

The eight layers documented here are the present state of that
carrying engineering. They are layered and validated. The geometry
itself is older than the engineering and will outlast it.


%% ==================================================================
\section*{Acknowledgments}
%% ==================================================================

Solo-authored. Drafting and editing of this manuscript were
assisted by Claude (Anthropic); the mathematical positions, design
choices, framing decisions, and acceptance or rejection of
empirical results are the author's. The Marcella runtime
referenced in \S\ref{sec:validation} is maintained by the Marcella
team at Davis Geometric; the cross-team A/B harness was run on the
public Marcella service against the live engine and is reproducible
from the artifacts cited in \S\ref{sec:validation}.

The DPU graphene processor referenced in
\S\ref{sec:validation-dpu} is the subject of U.S.\ Provisional
Application No.~64/073,981~\cite{davis2026dpu} (filed 2026-05-25).
The hardware realization of the substrate described in this paper is
currently in pre-fabrication design phase; physical device
characterization awaits a future companion publication.


%% ==================================================================
\begin{thebibliography}{99}

\bibitem{davis2026sheafcomposition}
B.~R.\ Davis,
\emph{Sheaf Composition: The Geometry of Creativity, Implemented},
Zenodo, 2026, doi:10.5281/zenodo.20185331.

\bibitem{davis2026purefiber}
B.~R.\ Davis,
\emph{Pure-Fiber Language Modeling},
Davis Geometric, May 2026.

\bibitem{davis2026doublecover}
B.~R.\ Davis,
\emph{The Double Cover Principle: Every Circle Carries Two
Circumferences},
Zenodo, 2026, doi:10.5281/zenodo.18895462.

\bibitem{davis2026manifold}
B.~R.\ Davis,
\emph{The Davis Manifold: A Riemannian Substrate for Residue-Bounded
Computation},
Davis Geometric, 2026. Zenodo DOI to be assigned.

\bibitem{gigi2026code}
B.~R.\ Davis,
\emph{GIGI: Geometric Intrinsic Global Index} (source code),
Davis Geometric, 2026.
Available at \url{https://github.com/nurdymuny/gigi/}.

\bibitem{davis2026davisduality}
B.~R.\ Davis,
\emph{The Davis Duality of Approximation and Obstruction: Why
Machine Learning Works, Why the Vacuum Has Mass, and the Universal
Law of Flat Failure},
Zenodo, 2026, doi:10.5281/zenodo.19428406.

\bibitem{davis2026dpu}
B.~R.\ Davis,
\emph{DPU System and Method for a Geometric Processor on a Graphene
Substrate via Berry-Phase Encoded Logic, K\"ahler-Topological
Protection, Multi-Channel Ballistic Transport, and Cross-Domain
Validated Prequantization Holonomy},
U.S.\ Provisional Application No.~64/073,981, filed May~25, 2026.

\bibitem{adachi2010hopfrinow}
T.~Adachi,
A theorem of Hopf--Rinow type for K\"ahler graphs,
\emph{Mem.\ Inst.\ Sci.\ Eng.\ Ritsumeikan Univ.}\ \textbf{69},
1--6 (2010).

\bibitem{bordemann1994toeplitz}
M.~Bordemann, E.~Meinrenken, M.~Schlichenmaier,
Toeplitz quantization of K\"ahler manifolds and $\mathrm{gl}(N)$,
$N \to \infty$ limits,
\emph{Comm.\ Math.\ Phys.}\ \textbf{165}, 281--296 (1994).

\bibitem{fhs2005chern}
T.~Fukui, Y.~Hatsugai, H.~Suzuki,
Chern Numbers in Discretized Brillouin Zone,
\emph{J.\ Phys.\ Soc.\ Jpn.}\ \textbf{74}, 1674--1677 (2005).

\bibitem{mccann2013bilayer}
E.~McCann, M.~Koshino,
The electronic properties of bilayer graphene,
\emph{Rep.\ Prog.\ Phys.}\ \textbf{76}, 056503 (2013).

\bibitem{wang2013contacts}
L.~Wang, I.~Meric, P.~Y.~Huang, et al.,
One-Dimensional Electrical Contact to a Two-Dimensional Material,
\emph{Science} \textbf{342}, 614--617 (2013).

\bibitem{berry1984quantal}
M.~V.\ Berry,
Quantal phase factors accompanying adiabatic changes,
\emph{Proc.\ Roy.\ Soc.\ A} \textbf{392}, 45--57 (1984).

\bibitem{arnold1989mechanics}
V.~I.\ Arnold,
\emph{Mathematical Methods of Classical Mechanics}, 2nd ed.,
Springer, 1989.

\bibitem{kobayashi1969foundations}
S.~Kobayashi and K.~Nomizu,
\emph{Foundations of Differential Geometry, Vols.~I and II},
Interscience, 1963 and 1969.

\bibitem{young2013pomdp}
S.~Young, M.~Ga\v{s}i\'c, B.~Thomson, and J.~D.\ Williams,
POMDP-based statistical spoken dialog systems: A review,
\emph{Proceedings of the IEEE} \textbf{101}(5), 1160--1179 (2013).

\bibitem{budzianowski2018multiwoz}
P.~Budzianowski, T.-H.\ Wen, B.-H.\ Tseng, I.~Casanueva,
S.~Ultes, O.~Ramadan, and M.~Ga\v{s}i\'c,
MultiWOZ -- A large-scale multi-domain Wizard-of-Oz dataset for
task-oriented dialogue modelling,
in \emph{Proc.\ EMNLP 2018}, pp.\ 5016--5026.

\bibitem{clarkbrennan1991grounding}
H.~H.\ Clark and S.~E.\ Brennan,
Grounding in communication, in
L.~B.\ Resnick, J.~M.\ Levine, and S.~D.\ Teasley (eds.),
\emph{Perspectives on Socially Shared Cognition},
American Psychological Association, 1991, pp.\ 127--149.

\end{thebibliography}

\end{document}
